% !TEX program = xelatex
\documentclass[UTF8,fontset=fandol,a4paper,11pt]{ctexart}
\usepackage[margin=21mm,top=23mm,bottom=22mm,headheight=16pt]{geometry}
\usepackage{amsmath,amssymb,mathtools,bm}
\usepackage{tikz}
\usetikzlibrary{arrows.meta,calc,angles,quotes,patterns}
\usepackage{xcolor,booktabs,array,longtable,enumitem,fancyhdr,needspace}
\usepackage[unicode,hidelinks]{hyperref}
\definecolor{ink}{HTML}{24343C}
\definecolor{accent}{HTML}{166A70}
\definecolor{muted}{HTML}{65747A}
\definecolor{rulegray}{HTML}{D4DDE0}
\definecolor{warm}{HTML}{966037}
\color{ink}
\setlength{\parindent}{2em}
\setlength{\parskip}{4pt}
\linespread{1.12}
\AtBeginDocument{\setlength{\abovedisplayskip}{6pt plus 2pt minus 1pt}
\setlength{\belowdisplayskip}{6pt plus 2pt minus 1pt}
\setlength{\abovedisplayshortskip}{0pt plus 1pt}
\setlength{\belowdisplayshortskip}{4pt plus 1pt}}
\setlist[enumerate]{label=(\arabic*),leftmargin=2.3em,itemsep=3pt,topsep=4pt}
\setlist[itemize]{leftmargin=2em,itemsep=2pt,topsep=4pt}
\allowdisplaybreaks[1]
\setlength{\emergencystretch}{2em}
\setcounter{secnumdepth}{0}
\providecommand{\tightlist}{\setlength{\itemsep}{0pt}\setlength{\parskip}{0pt}}
\ctexset{section={format=\Large\bfseries\color{accent},beforeskip=10pt,afterskip=9pt},
subsection={format=\normalsize\bfseries\color{accent},beforeskip=10pt,afterskip=4pt},
subsubsection={format=\normalsize\bfseries,beforeskip=8pt,afterskip=3pt}}
\pagestyle{fancy}
\fancyhf{}
\fancyhead[L]{\small\color{muted}物理学难题集萃}
% 页眉仅保留书名；版本名称只在封面显示。
\fancyfoot[L]{\small\color{muted}\chapterlabel}
\fancyfoot[R]{\small\thepage}
\renewcommand{\headrulewidth}{0.4pt}
\renewcommand{\footrulewidth}{0pt}
\newcommand{\dd}{\mathop{}\!\mathrm{d}}
\newcommand{\vect}[1]{\bm{#1}}
\newcommand{\tagline}[1]{\noindent{\small\color{muted}#1}\par}
\newcommand{\result}[1]{\begin{equation*}\boxed{#1}\end{equation*}}
\newcommand{\figcaption}[1]{\par\vspace{2pt}{\small\color{muted}#1}\par}
\newcommand{\teach}[2]{\Needspace{4\baselineskip}\subsection*{讲解要点与追问}#1\par\smallskip\noindent\textbf{追问：}#2\par}
\newcommand{\workarea}{\par\vspace{7mm}\noindent{\small\color{muted}作答区}\par\vspace{4mm}
\noindent\textcolor{rulegray}{\rule{\linewidth}{0.3pt}}\par\vspace{16mm}
\noindent\textcolor{rulegray}{\rule{\linewidth}{0.3pt}}\par\vspace{16mm}
\noindent\textcolor{rulegray}{\rule{\linewidth}{0.3pt}}\par\vspace{16mm}
\noindent\textcolor{rulegray}{\rule{\linewidth}{0.3pt}}}
\tikzset{axis/.style={-{Stealth[length=2mm]},thin,color=muted},
motion/.style={-{Stealth[length=2.3mm]},thick,color=accent},
force/.style={-{Stealth[length=2.3mm]},thick,color=warm},
traj/.style={thick,color=accent},guide/.style={dashed,thin,color=muted},
object/.style={thick,color=ink},every node/.style={font=\small}}\widowpenalty=10000\clubpenalty=10000\displaywidowpenalty=10000\newcommand{\chapterlabel}{力学试题（二）}







\hypersetup{pdftitle={力学试题二 力学试题（二）·学生解答版},pdfauthor={}}

\begin{document}\pdfbookmark[0]{力学试题二 力学试题（二）}{cover}

\thispagestyle{empty}\vspace*{22mm}

\noindent{\Huge\bfseries 物理学难题集萃}\par\vspace{10mm}

\noindent{\LARGE 力学试题二}\par\vspace{6mm}

\noindent{\LARGE 力学试题（二）}\par\vspace{12mm}

\noindent{\Large\color{accent}学生解答版}\par\vspace{9mm}

\noindent\textcolor{accent}{\rule{\linewidth}{1pt}}\par\vspace{9mm}

\noindent 第1--6题\par\vfill\clearpage

\pdfbookmark[1]{题目目录}{contents}\section*{题目目录}

\begin{longtable}{@{}r p{113mm} r@{}}\toprule 题号 & 题目 & 页码\\\midrule

\hyperref[q:1]{1} & \hyperref[q:1]{移动圆环与带环套的地面轨迹} & \pageref{q:1} \\

\hyperref[q:2]{2} & \hyperref[q:2]{三次喷气中的能量、速度矢量与落地点} & \pageref{q:2} \\

\hyperref[q:3]{3} & \hyperref[q:3]{绳越过凸圆弧时的张力与脱离} & \pageref{q:3} \\

\hyperref[q:4]{4} & \hyperref[q:4]{两细杆端部弹性碰撞与接触点速度} & \pageref{q:4} \\

\hyperref[q:5]{5} & \hyperref[q:5]{对称铰接双杆剪绳后的下落} & \pageref{q:5} \\

\hyperref[q:6]{6} & \hyperref[q:6]{双线悬圆柱的平移、复摆与偏航} & \pageref{q:6} \\

\bottomrule\end{longtable}

\clearpage

\pdfbookmark[2]{第 1 题 移动圆环与带环套的地面轨迹}{q1}

\section*{第 1 题\quad 移动圆环与带环套的地面轨迹}\label{q:1}


光滑水平面上放着均匀圆环，质量为$M$、半径为$R$。质量为$m<M$的质点可沿环的光滑内壁滑动。开始时圆环静止，环心在$O=(0,0)$，质点位于$(0,R)$，初速度为$v_0\hat x$。运动始终在水平面内进行。
\begin{enumerate}
\item 证明质点不会离开内壁。
\item 求质点相对地面的运动方程，并画出轨迹。
\item 求轨迹在最高处、最低处和竖直切线处的曲率半径；说明一周期内各类点出现的位置。
\end{enumerate}


\begin{center}\begin{minipage}{.98\linewidth}\centering
\begin{tikzpicture}[scale=.9,transform shape]
\draw[axis] (0.00000,0.00000)--(4.25000,0.00000);\draw[axis] (0.00000,-3.50000)--(0.00000,4.25000);\node[below right] at (4.25000,0.00000) {$x$};\node[above right] at (0.00000,4.25000) {$y$};\draw[thin,color=muted] (-3.50000,0.00000)--(0.00000,0.00000);\draw[object] (0.00000,0.00000) circle(3.00000);\draw[object] (0.00000,0.00000) circle(3.10000);\fill (0.00000,3.00000) circle(.035);\node[above left] at (0.00000,3.00000) {$m$};\fill (0.00000,0.00000) circle(.035);\node[below left] at (0.00000,0.00000) {$O$};\draw[motion] (0.00000,3.00000)--(2.00000,3.00000);\node[above] at (1.00000,3.10000) {$v_0$};\draw[<->,thin,color=muted] (0.00000,-0.60000)--(3.00000,-0.60000);\node[below] at (1.50000,-0.60000) {$R$};\node[below] at (-2.00000,-3.60000) {圆环 $M$，水平面};
\end{tikzpicture}
\figcaption{图力试2.1\quad 初始质点在环内壁顶端，初速沿+x，运动在水平面内。}
\end{minipage}\end{center}

\Needspace{6\baselineskip}\subsection*{解答}


图力试2.1中，圆环与质点的水平总动量初始为$mv_0\hat x$。系统质心$C$的速度和初始纵坐标为
\[
\boldsymbol V_C=\frac{m}{M+m}v_0\hat x,\qquad
y_C(0)=b=\frac{mR}{M+m}.
\]
质心系是匀速惯性系。记
\[
a=\frac{MR}{M+m},\qquad b=\frac{mR}{M+m},\qquad a+b=R.
\]
质点与环心分别位于$C$两侧，距离为$a,b$。内壁作用力通过环心，不产生圆环自转的力矩；环心仍会在该力下平移。
这里把两种角运动分开：圆环相对自身环心的转动满足$I_O\dot\Omega_{\rm body}=0$，初始自身角速率为零，故始终保持零；环心在质心系的运动则由水平合力决定。水平面内的相互作用改变各自的速度，系统总动量保持。


\begin{center}\begin{minipage}{.98\linewidth}\centering
\begin{tikzpicture}[scale=.76,transform shape]
\draw[object] (-1.00000,0.00000) circle(3.00000);\draw[guide] (0.00000,0.00000) circle(2.00000);\draw[guide] (0.00000,0.00000) circle(1.00000);\fill (-1.00000,0.00000) circle(.035);\node[below left] at (-1.00000,0.00000) {环心};\fill (0.00000,0.00000) circle(.035);\node[above] at (0.00000,0.00000) {$C$};\fill (2.00000,0.00000) circle(.035);\node[below right] at (2.00000,0.00000) {$m$};\draw[motion] (2.00000,0.00000)--(2.00000,-1.50000);\draw[motion] (-1.00000,0.00000)--(-1.00000,1.50000);\node[right] at (2.10000,-1.20000) {$a\omega$};\node[left] at (-1.10000,1.20000) {$b\omega$};\draw[force] (2.00000,0.00000)--(0.80000,0.00000);\node[above] at (1.50000,0.10000) {$N$};\draw[<->,thin,color=muted] (-1.00000,-3.60000)--(0.00000,-3.60000);\node[below] at (-0.50000,-3.60000) {$b$};\draw[<->,thin,color=muted] (0.00000,-4.50000)--(2.00000,-4.50000);\node[below] at (1.00000,-4.50000) {$a$};\draw[motion] (-1.00000,4.00000)--(1.00000,4.00000);\node[above] at (0.00000,4.10000) {质心系以 $V_C$向右移动};
\end{tikzpicture}
\figcaption{解图力试2.1A\quad 质心系中两圆半径a、b，内壁支持力始终向内。}
\end{minipage}\end{center}
（1）解图力试2.1A先在质心系看接触运动。相互作用力沿环心与质点的连线，也通过$C$，所以质点关于$C$的角动量守恒。保持接触时半径$a$不变，故其角速率为常量。初始相对速度为
\[
v'_m=v_0-V_C=\frac{M}{M+m}v_0,\qquad
\omega=\frac{v'_m}{a}=\frac{v_0}{R}.
\]
环心的初始质心系速度为$v'_O=-V_C=-mv_0/(M+m)$，其大小除以$b$同样得到$v_0/R$。质点在$C$上方、环心在$C$下方，两速度方向相反而共同绕$C$顺时针运动。两个物体的径向方程也可合在一起：若$\boldsymbol n$从环心指向质点，则相对加速度为$\ddot{\boldsymbol r}=-N(1/m+1/M)\boldsymbol n$，接触时$|\boldsymbol r|=R$给$N(1/m+1/M)=R\omega^2$。这个正向心需求与角动量确定的匀速圆周相接。
圆环内壁对质点所需支持力为
\[
\boxed{N=ma\omega^2=\frac{Mm}{M+m}\frac{v_0^2}{R}>0}.
\]
该力始终指向环心，由内壁对质点的正压力提供。保持接触的圆周解在每一相位都要求同样的正支持力，满足内壁单侧接触条件，因而该接触解持续成立。

接触解中的质点、环心始终分居质心两侧，且$ma=Mb$。取顺时针径向单位矢量$\boldsymbol n=(\sin\tau,\cos\tau)$，两位置相对$C$分别为$a\boldsymbol n$与$-b\boldsymbol n$，相减为$R\boldsymbol n$。因此环心虽然移动，圆环的几何约束半径仍严格等于$R$。质心平移再加到这两位置上时，对两物体加的是同一个$V_Ct\hat x$，保持两者的相对运动与接触力。
（2）令$\tau=\omega t$。质点在质心系中的位置为$(a\sin\tau,a\cos\tau)$，质心在地面系中为$(V_Ct,b)$，于是
\[
\boxed{x=V_Ct+a\sin\omega t,\qquad
y=b+a\cos\omega t}.
\]
用$V_C=b\omega$，也可写成$x=b\tau+a\sin\tau$。环心则为
\[
X=V_Ct-b\sin\tau,\qquad Y=b-b\cos\tau.
\]
环心速度从自己的参数式得到$\dot X=b\omega(1-\cos\tau)$、$\dot Y=b\omega\sin\tau$，出发时两分量均为零，符合环初始静止。质点在出发时的水平速度为$(a+b)\omega=v_0$，两初始状态因此由同一平移和相反圆周运动共同实现。
初始条件$x(0)=0,y(0)=R,\dot x(0)=v_0,\dot y(0)=0$均满足。解图力试2.1B中$a>b$，圆周运动叠加较慢的匀速平移，轨迹产生环套；一圈以后整体沿$x$平移$2\pi b$。
\begin{center}\begin{minipage}{.98\linewidth}\centering
\begin{tikzpicture}[scale=.72,transform shape]
\draw[axis] (0.00000,0.00000)--(13.25000,0.00000);\draw[axis] (0.00000,-1.80000)--(0.00000,4.05000);\node[below right] at (13.25000,0.00000) {$x$};\node[above right] at (0.00000,4.05000) {$y$};\draw[traj] plot coordinates {(0.00000,3.00000) (0.07554,2.99937) (0.15106,2.99746) (0.22650,2.99430) (0.30186,2.98986) (0.37708,2.98417) (0.45215,2.97721) (0.52702,2.96901) (0.60167,2.95955) (0.67607,2.94885) (0.75019,2.93692) (0.82398,2.92375) (0.89743,2.90937) (0.97051,2.89378) (1.04317,2.87698) (1.11540,2.85900) (1.18716,2.83983) (1.25842,2.81950) (1.32916,2.79802) (1.39934,2.77539) (1.46893,2.75164) (1.53792,2.72678) (1.60626,2.70082) (1.67394,2.67379) (1.74092,2.64569) (1.80718,2.61655) (1.87270,2.58639) (1.93744,2.55522) (2.00139,2.52306) (2.06451,2.48994) (2.12679,2.45587) (2.18820,2.42088) (2.24871,2.38498) (2.30831,2.34821) (2.36697,2.31059) (2.42468,2.27213) (2.48141,2.23287) (2.53713,2.19282) (2.59184,2.15202) (2.64552,2.11049) (2.69814,2.06825) (2.74968,2.02534) (2.80014,1.98177) (2.84949,1.93758) (2.89772,1.89280) (2.94482,1.84745) (2.99077,1.80157) (3.03555,1.75518) (3.07916,1.70830) (3.12159,1.66098) (3.16282,1.61324) (3.20284,1.56511) (3.24164,1.51662) (3.27922,1.46781) (3.31557,1.41870) (3.35068,1.36932) (3.38454,1.31971) (3.41714,1.26989) (3.44849,1.21991) (3.47858,1.16978) (3.50741,1.11955) (3.53497,1.06924) (3.56126,1.01889) (3.58629,0.96852) (3.61004,0.91818) (3.63253,0.86788) (3.65376,0.81768) (3.67372,0.76758) (3.69242,0.71764) (3.70986,0.66787) (3.72606,0.61831) (3.74101,0.56900) (3.75472,0.51996) (3.76720,0.47122) (3.77846,0.42282) (3.78850,0.37479) (3.79734,0.32715) (3.80498,0.27993) (3.81144,0.23318) (3.81673,0.18691) (3.82086,0.14116) (3.82384,0.09595) (3.82569,0.05131) (3.82643,0.00728) (3.82607,-0.03613) (3.82462,-0.07887) (3.82210,-0.12094) (3.81853,-0.16229) (3.81393,-0.20290) (3.80832,-0.24276) (3.80171,-0.28182) (3.79413,-0.32007) (3.78560,-0.35749) (3.77613,-0.39404) (3.76576,-0.42971) (3.75450,-0.46447) (3.74237,-0.49831) (3.72941,-0.53119) (3.71563,-0.56310) (3.70106,-0.59402) (3.68572,-0.62393) (3.66964,-0.65282) (3.65285,-0.68065) (3.63537,-0.70741) (3.61723,-0.73310) (3.59845,-0.75768) (3.57907,-0.78115) (3.55912,-0.80350) (3.53861,-0.82469) (3.51759,-0.84473) (3.49608,-0.86360) (3.47410,-0.88129) (3.45170,-0.89779) (3.42890,-0.91308) (3.40572,-0.92716) (3.38221,-0.94001) (3.35839,-0.95164) (3.33429,-0.96203) (3.30994,-0.97117) (3.28538,-0.97907) (3.26064,-0.98571) (3.23575,-0.99109) (3.21074,-0.99521) (3.18563,-0.99806) (3.16048,-0.99964) (3.13530,-0.99996) (3.11012,-0.99901) (3.08499,-0.99679) (3.05993,-0.99331) (3.03497,-0.98856) (3.01015,-0.98255) (2.98550,-0.97528) (2.96104,-0.96676) (2.93681,-0.95699) (2.91285,-0.94598) (2.88918,-0.93374) (2.86583,-0.92027) (2.84284,-0.90558) (2.82023,-0.88969) (2.79804,-0.87260) (2.77629,-0.85432) (2.75502,-0.83486) (2.73425,-0.81424) (2.71402,-0.79247) (2.69435,-0.76956) (2.67527,-0.74553) (2.65681,-0.72039) (2.63900,-0.69417) (2.62186,-0.66686) (2.60542,-0.63850) (2.58971,-0.60911) (2.57475,-0.57869) (2.56057,-0.54727) (2.54719,-0.51487) (2.53464,-0.48151) (2.52295,-0.44720) (2.51213,-0.41199) (2.50220,-0.37587) (2.49320,-0.33888) (2.48514,-0.30105) (2.47805,-0.26239) (2.47193,-0.22293) (2.46683,-0.18269) (2.46274,-0.14170) (2.45969,-0.09999) (2.45771,-0.05758) (2.45680,-0.01450) (2.45698,0.02922) (2.45828,0.07356) (2.46069,0.11848) (2.46425,0.16397) (2.46896,0.20998) (2.47483,0.25650) (2.48188,0.30349) (2.49012,0.35091) (2.49956,0.39876) (2.51021,0.44698) (2.52207,0.49555) (2.53517,0.54444) (2.54950,0.59362) (2.56507,0.64306) (2.58189,0.69273) (2.59996,0.74259) (2.61929,0.79261) (2.63988,0.84277) (2.66174,0.89302) (2.68486,0.94335) (2.70925,0.99370) (2.73491,1.04407) (2.76184,1.09440) (2.79003,1.14468) (2.81949,1.19486) (2.85021,1.24492) (2.88219,1.29482) (2.91542,1.34454) (2.94991,1.39404) (2.98563,1.44329) (3.02260,1.49226) (3.06079,1.54091) (3.10020,1.58922) (3.14083,1.63716) (3.18266,1.68470) (3.22568,1.73180) (3.26988,1.77843) (3.31525,1.82458) (3.36177,1.87020) (3.40944,1.91527) (3.45823,1.95975) (3.50814,2.00363) (3.55914,2.04688) (3.61123,2.08945) (3.66437,2.13134) (3.71857,2.17251) (3.77379,2.21294) (3.83002,2.25260) (3.88724,2.29146) (3.94542,2.32950) (4.00456,2.36671) (4.06462,2.40304) (4.12558,2.43849) (4.18743,2.47302) (4.25013,2.50662) (4.31367,2.53926) (4.37802,2.57093) (4.44315,2.60160) (4.50904,2.63125) (4.57567,2.65987) (4.64300,2.68744) (4.71101,2.71394) (4.77968,2.73935) (4.84897,2.76366) (4.91887,2.78685) (4.98933,2.80890) (5.06033,2.82981) (5.13184,2.84956) (5.20384,2.86814) (5.27629,2.88553) (5.34917,2.90172) (5.42243,2.91671) (5.49606,2.93049) (5.57002,2.94304) (5.64428,2.95435) (5.71881,2.96443) (5.79357,2.97327) (5.86855,2.98085) (5.94370,2.98717) (6.01899,2.99224) (6.09440,2.99604) (6.16988,2.99857) (6.24541,2.99984) (6.32096,2.99984) (6.39649,2.99857) (6.47198,2.99604) (6.54738,2.99224) (6.62267,2.98717) (6.69782,2.98085) (6.77280,2.97327) (6.84756,2.96443) (6.92209,2.95435) (6.99635,2.94304) (7.07031,2.93049) (7.14394,2.91671) (7.21721,2.90172) (7.29008,2.88553) (7.36253,2.86814) (7.43453,2.84956) (7.50604,2.82981) (7.57704,2.80890) (7.64750,2.78685) (7.71740,2.76366) (7.78669,2.73935) (7.85536,2.71394) (7.92337,2.68744) (7.99070,2.65987) (8.05733,2.63125) (8.12322,2.60160) (8.18835,2.57093) (8.25270,2.53926) (8.31624,2.50662) (8.37894,2.47302) (8.44079,2.43849) (8.50175,2.40304) (8.56181,2.36671) (8.62095,2.32950) (8.67913,2.29146) (8.73635,2.25260) (8.79258,2.21294) (8.84780,2.17251) (8.90200,2.13134) (8.95514,2.08945) (9.00723,2.04688) (9.05823,2.00363) (9.10814,1.95975) (9.15693,1.91527) (9.20460,1.87020) (9.25112,1.82458) (9.29649,1.77843) (9.34069,1.73180) (9.38371,1.68470) (9.42554,1.63716) (9.46617,1.58922) (9.50558,1.54091) (9.54377,1.49226) (9.58074,1.44329) (9.61646,1.39404) (9.65095,1.34454) (9.68418,1.29482) (9.71616,1.24492) (9.74688,1.19486) (9.77634,1.14468) (9.80453,1.09440) (9.83146,1.04407) (9.85712,0.99370) (9.88151,0.94335) (9.90463,0.89302) (9.92649,0.84277) (9.94708,0.79261) (9.96641,0.74259) (9.98448,0.69273) (10.00130,0.64306) (10.01687,0.59362) (10.03120,0.54444) (10.04430,0.49555) (10.05616,0.44698) (10.06681,0.39876) (10.07625,0.35091) (10.08449,0.30349) (10.09154,0.25650) (10.09741,0.20998) (10.10212,0.16397) (10.10568,0.11848) (10.10809,0.07356) (10.10939,0.02922) (10.10957,-0.01450) (10.10866,-0.05758) (10.10668,-0.09999) (10.10363,-0.14170) (10.09954,-0.18269) (10.09444,-0.22293) (10.08832,-0.26239) (10.08123,-0.30105) (10.07317,-0.33888) (10.06417,-0.37587) (10.05424,-0.41199) (10.04342,-0.44720) (10.03173,-0.48151) (10.01918,-0.51487) (10.00580,-0.54727) (9.99162,-0.57869) (9.97667,-0.60911) (9.96095,-0.63850) (9.94452,-0.66686) (9.92738,-0.69417) (9.90956,-0.72039) (9.89110,-0.74553) (9.87202,-0.76956) (9.85235,-0.79247) (9.83212,-0.81424) (9.81135,-0.83486) (9.79008,-0.85432) (9.76833,-0.87260) (9.74614,-0.88969) (9.72353,-0.90558) (9.70054,-0.92027) (9.67719,-0.93374) (9.65352,-0.94598) (9.62956,-0.95699) (9.60533,-0.96676) (9.58087,-0.97528) (9.55622,-0.98255) (9.53140,-0.98856) (9.50644,-0.99331) (9.48138,-0.99679) (9.45625,-0.99901) (9.43107,-0.99996) (9.40589,-0.99964) (9.38074,-0.99806) (9.35564,-0.99521) (9.33062,-0.99109) (9.30573,-0.98571) (9.28099,-0.97907) (9.25643,-0.97117) (9.23208,-0.96203) (9.20798,-0.95164) (9.18416,-0.94001) (9.16065,-0.92716) (9.13747,-0.91308) (9.11467,-0.89779) (9.09227,-0.88129) (9.07029,-0.86360) (9.04878,-0.84473) (9.02776,-0.82469) (9.00725,-0.80350) (8.98730,-0.78115) (8.96792,-0.75768) (8.94915,-0.73310) (8.93100,-0.70741) (8.91352,-0.68065) (8.89673,-0.65282) (8.88065,-0.62393) (8.86532,-0.59402) (8.85074,-0.56310) (8.83696,-0.53119) (8.82400,-0.49831) (8.81187,-0.46447) (8.80061,-0.42971) (8.79024,-0.39404) (8.78077,-0.35749) (8.77224,-0.32007) (8.76466,-0.28182) (8.75805,-0.24276) (8.75244,-0.20290) (8.74784,-0.16229) (8.74427,-0.12094) (8.74175,-0.07887) (8.74030,-0.03613) (8.73994,0.00728) (8.74068,0.05131) (8.74253,0.09595) (8.74551,0.14116) (8.74964,0.18691) (8.75493,0.23318) (8.76139,0.27993) (8.76904,0.32715) (8.77787,0.37479) (8.78791,0.42282) (8.79917,0.47122) (8.81165,0.51996) (8.82536,0.56900) (8.84031,0.61831) (8.85651,0.66787) (8.87395,0.71764) (8.89265,0.76758) (8.91261,0.81768) (8.93384,0.86788) (8.95633,0.91818) (8.98008,0.96852) (9.00511,1.01889) (9.03140,1.06924) (9.05896,1.11955) (9.08779,1.16978) (9.11788,1.21991) (9.14923,1.26989) (9.18184,1.31971) (9.21569,1.36932) (9.25080,1.41870) (9.28715,1.46781) (9.32473,1.51662) (9.36353,1.56511) (9.40355,1.61324) (9.44478,1.66098) (9.48721,1.70830) (9.53082,1.75518) (9.57560,1.80157) (9.62155,1.84745) (9.66865,1.89280) (9.71688,1.93758) (9.76623,1.98177) (9.81669,2.02534) (9.86823,2.06825) (9.92085,2.11049) (9.97453,2.15202) (10.02924,2.19282) (10.08496,2.23287) (10.14169,2.27213) (10.19940,2.31059) (10.25806,2.34821) (10.31766,2.38498) (10.37817,2.42088) (10.43958,2.45587) (10.50186,2.48994) (10.56498,2.52306) (10.62893,2.55522) (10.69367,2.58639) (10.75919,2.61655) (10.82545,2.64569) (10.89243,2.67379) (10.96011,2.70082) (11.02845,2.72678) (11.09744,2.75164) (11.16703,2.77539) (11.23721,2.79802) (11.30795,2.81950) (11.37921,2.83983) (11.45097,2.85900) (11.52320,2.87698) (11.59586,2.89378) (11.66894,2.90937) (11.74239,2.92375) (11.81618,2.93692) (11.89030,2.94885) (11.96470,2.95955) (12.03935,2.96901) (12.11422,2.97721) (12.18929,2.98417) (12.26451,2.98986) (12.33987,2.99430) (12.41531,2.99746) (12.49083,2.99937) (12.56637,3.00000)};\fill (0.00000,3.00000) circle(.035);\node[above left] at (0.00000,3.00000) {$A_1$};\fill (3.14159,-1.00000) circle(.035);\node[below right] at (3.14159,-1.00000) {$A_2$};\fill (6.28319,3.00000) circle(.035);\node[above] at (6.28319,3.00000) {$A_1$};\fill (3.82645,0.00000) circle(.035);\node[above right] at (3.82645,0.00000) {$A_{3a}$};\fill (2.45674,-0.00000) circle(.035);\node[below left] at (2.45674,-0.00000) {$A_{3b}$};\draw[<->,thin,color=muted] (0.00000,-2.70000)--(6.28319,-2.70000);\node[below] at (3.14159,-2.70000) {$2\pi b$};\node[above] at (6.00000,4.50000) {示例 $M=2m$；图示长度单位一致};
\end{tikzpicture}
\figcaption{解图力试2.1B\quad 匀速平移叠加圆运动，产生环套与三类特殊切线点。}
\end{minipage}\end{center}

（3）对参数方程求导：
\[
v_x=\frac{v_0}{M+m}(m+M\cos\tau),\qquad
v_y=-\frac{Mv_0}{M+m}\sin\tau,
\]
\[
a_x=-\frac{Mv_0^2}{R(M+m)}\sin\tau,\qquad
a_y=-\frac{Mv_0^2}{R(M+m)}\cos\tau.
\]
总加速度大小为常量
\[
A_0=\frac{Mv_0^2}{R(M+m)}.
\]
一般位置的加速度同时含切向与法向分量，法向部分按地面速度的垂直方向投影。

\begin{center}\begin{minipage}{.98\linewidth}\centering
\begin{tikzpicture}[scale=.64,transform shape]
\fill (0.00000,0.00000) circle(.035);\node[below left] at (0.00000,0.00000) {$A_1$};\draw[motion] (0.15000,0.15000)--(2.15000,0.15000);\node[above] at (2.00000,0.60000) {$\boldsymbol v$};\draw[force] (0.00000,0.00000)--(0.00000,-1.50000);\node[below left] at (0.00000,-1.50000) {$a_n$};\begin{scope}[shift={(7.00000,0.00000)}]
\fill (0.00000,0.00000) circle(.035);\node[below left] at (0.00000,0.00000) {$A_2$};\draw[motion] (0.15000,0.15000)--(-1.85000,0.15000);\node[above] at (-2.00000,0.60000) {$\boldsymbol v$};\draw[force] (0.00000,0.00000)--(0.00000,1.50000);\node[below left] at (0.00000,1.50000) {$a_n$};
\end{scope}\begin{scope}[shift={(14.00000,0.00000)}]
\fill (0.00000,0.00000) circle(.035);\node[above right] at (0.00000,0.00000) {$A_3$};\draw[motion] (0.15000,0.15000)--(0.15000,-1.85000);\node[right] at (0.25000,-2.10000) {$\boldsymbol v$};\draw[force] (0.00000,0.00000)--(-1.50000,0.00000);\node[below left] at (-1.50000,0.00000) {$a_n$};\draw[guide] (0.00000,0.00000)--(0.00000,0.90000);\node[above left] at (0.00000,0.90000) {$a_t$};
\end{scope}\node[below] at (7.00000,-3.20000) {竖直切线处只取水平法向分量};
\end{tikzpicture}
\figcaption{解图力试2.1C\quad 最高、最低点加速度全为法向；竖切点还含切向分量。}
\end{minipage}\end{center}
解图力试2.1C分别画三个代表点的速度、法向加速度。最高点$A_1$满足$\tau=2k\pi$，速度沿$+x$且$v=v_0$，加速度竖直向下，故
\[
\boxed{\rho_1=\frac{v_0^2}{A_0}=\frac{M+m}{M}R}.
\]
最低点$A_2$满足$\tau=(2k+1)\pi$，其地面速度向左，速率为$(M-m)v_0/(M+m)$，加速度竖直向上，所以
\[
\boxed{\rho_2=\frac{(M-m)^2}{M(M+m)}R}.
\]
向左运动的符号决定切线方向，曲率半径使用速率平方和法向加速度大小。

竖直切线点$A_3$由$v_x=0$确定，即
\[
\cos\tau=-\frac mM.
\]
由于$m<M$，每周期有两个这样的点。
记$\tau_3=\arccos(-m/M)$，一个周期的两点相位为$\tau_3$和$2\pi-\tau_3$。前者$\sin\tau>0$，质点速度向下、水平法向加速度向左；后者两方向分别向上、向右。两点的竖直加速度同为$A_0m/M$，它属于沿切线的分量，所以两点都只取水平分量计算法向加速度。
此时
\[
v^2=v_y^2=\frac{M-m}{M+m}v_0^2,\qquad
a_n=|a_x|=\frac{v_0^2}{R}\sqrt{\frac{M-m}{M+m}},
\]
因法线水平而取$|a_x|$，得到
\[
\boxed{\rho_3=R\sqrt{\frac{M-m}{M+m}}}.
\]
各类点的位置也可直接标在地面图上。最高点为$(2k\pi b,R)$，最低点为$((2k+1)\pi b,b-a)$。竖直切线处$y=b-a(m/M)=0$，横坐标分别为$b\tau_3+a\sqrt{1-m^2/M^2}$和$b(2\pi-\tau_3)-a\sqrt{1-m^2/M^2}$，再按周期整体平移$2k\pi b$。这些坐标将三类曲率与实际轨道上的点逐一对应。

一般点的速率满足$v^2=\omega^2(a^2+b^2+2ab\cos\tau)$，速度与加速度的平面叉积为$v_xa_y-v_ya_x=-a\omega^3(a+b\cos\tau)$。其绝对值除以速率就是法向加速度，随后用$v^2/a_n$可统一得到各点曲率半径。$a>b$时$a+b\cos\tau\ge a-b>0$，所以全过程曲线光滑，法向与运动方向按同一连续支变化。
同一局部法向关系也可由参数式统一写为
\[
\rho(\tau)=
\frac{(a^2+b^2+2ab\cos\tau)^{3/2}}
{a(a+b\cos\tau)}.
\]
本题$a>b$使分母恒正；在上述三类点代入，分别恢复三个结果。曲线交叉的两次经过属于不同的$\tau$，各次经过分别使用其对应时刻的速度切线，两支保持各自的参数连续性。
\begin{center}\begin{minipage}{.98\linewidth}\centering
\begin{tikzpicture}[scale=.8,transform shape]
\draw[axis] (0,0)--(10.7,0) node[right] {$\tau$};
\draw[axis] (0,0)--(0,3.2) node[above] {$\rho/R$};
\node[above] at (5,3.5) {$M=2m$时一周期的地面曲率半径};
\draw[traj] plot coordinates {(0.000000000,2.700000000) (0.083333333,2.698766393) (0.166666667,2.695066702) (0.250000000,2.688904314) (0.333333333,2.680284894) (0.416666667,2.669216401) (0.500000000,2.655709126) (0.583333333,2.639775729) (0.666666667,2.621431296) (0.750000000,2.600693403) (0.833333333,2.577582189) (0.916666667,2.552120449) (1.000000000,2.524333738) (1.083333333,2.494250491) (1.166666667,2.461902160) (1.250000000,2.427323372) (1.333333333,2.390552104) (1.416666667,2.351629888) (1.500000000,2.310602034) (1.583333333,2.267517886) (1.666666667,2.222431101) (1.750000000,2.175399970) (1.833333333,2.126487765) (1.916666667,2.075763126) (2.000000000,2.023300488) (2.083333333,1.969180547) (2.166666667,1.913490765) (2.250000000,1.856325916) (2.333333333,1.797788678) (2.416666667,1.737990245) (2.500000000,1.677050983) (2.583333333,1.615101097) (2.666666667,1.552281307) (2.750000000,1.488743516) (2.833333333,1.424651437) (2.916666667,1.360181165) (3.000000000,1.295521631) (3.083333333,1.230874929) (3.166666667,1.166456418) (3.250000000,1.102494572) (3.333333333,1.039230485) (3.416666667,0.976916967) (3.500000000,0.915817149) (3.583333333,0.856202529) (3.666666667,0.798350396) (3.750000000,0.742540609) (3.833333333,0.689051736) (3.916666667,0.638156612) (4.000000000,0.590117450) (4.083333333,0.545180719) (4.166666667,0.503572094) (4.250000000,0.465491853) (4.333333333,0.431111170) (4.416666667,0.400569746) (4.500000000,0.373975188) (4.583333333,0.351404386) (4.666666667,0.332906924) (4.750000000,0.318510237) (4.833333333,0.308225916) (4.916666667,0.302056257) (5.000000000,0.300000000) (5.083333333,0.302056257) (5.166666667,0.308225916) (5.250000000,0.318510237) (5.333333333,0.332906924) (5.416666667,0.351404386) (5.500000000,0.373975188) (5.583333333,0.400569746) (5.666666667,0.431111170) (5.750000000,0.465491853) (5.833333333,0.503572094) (5.916666667,0.545180719) (6.000000000,0.590117450) (6.083333333,0.638156612) (6.166666667,0.689051736) (6.250000000,0.742540609) (6.333333333,0.798350396) (6.416666667,0.856202529) (6.500000000,0.915817149) (6.583333333,0.976916967) (6.666666667,1.039230485) (6.750000000,1.102494572) (6.833333333,1.166456418) (6.916666667,1.230874929) (7.000000000,1.295521631) (7.083333333,1.360181165) (7.166666667,1.424651437) (7.250000000,1.488743516) (7.333333333,1.552281307) (7.416666667,1.615101097) (7.500000000,1.677050983) (7.583333333,1.737990245) (7.666666667,1.797788678) (7.750000000,1.856325916) (7.833333333,1.913490765) (7.916666667,1.969180547) (8.000000000,2.023300488) (8.083333333,2.075763126) (8.166666667,2.126487765) (8.250000000,2.175399970) (8.333333333,2.222431101) (8.416666667,2.267517886) (8.500000000,2.310602034) (8.583333333,2.351629888) (8.666666667,2.390552104) (8.750000000,2.427323372) (8.833333333,2.461902160) (8.916666667,2.494250491) (9.000000000,2.524333738) (9.083333333,2.552120449) (9.166666667,2.577582189) (9.250000000,2.600693403) (9.333333333,2.621431296) (9.416666667,2.639775729) (9.500000000,2.655709126) (9.583333333,2.669216401) (9.666666667,2.680284894) (9.750000000,2.688904314) (9.833333333,2.695066702) (9.916666667,2.698766393) (10.000000000,2.700000000)};
\node[below] at (0,-.12) {$0$};
\node[below] at (5,-.12) {$\pi$};
\node[below] at (10,-.12) {$2\pi$};
\end{tikzpicture}
\figcaption{解图力试2.1D\quad $M=2m$时一周期的地面曲率半径}
\end{minipage}\end{center}
图D中最高点的曲率半径最大，最低点最小；两竖直切线点落在相同的中间曲率值。每一时刻都用地面速度确定法线，质心系的匀速圆周则提供保持不变的加速度矢量。


\Needspace{13\baselineskip}

\pdfbookmark[2]{第 2 题 三次喷气中的能量、速度矢量与落地点}{q2}

\section*{第 2 题\quad 三次喷气中的能量、速度矢量与落地点}\label{q:2}


飞船起初在半径$r_0$的圆轨道上绕地球运动。每次喷气持续时间极短，喷气过程中位置不变，各阶段把飞船质量视为常量。将单位质量的冲量$\Delta\boldsymbol v=\Delta\boldsymbol p/m$称为比冲量，其大小单位为速度。地球质量为$M_E$、半径为$R_E$，$r_1>r_0>R_E$，忽略空气阻力。
\begin{enumerate}
\item 第一次喷气使飞船刚能逃逸，求最小比冲量及排气方向。
\item 第一次取上述最小值。飞船沿上升轨迹到达$r_1$处，第二次喷气使它进入半径$r_1$的同向圆轨道。求比冲量及排气方向。
\item 飞船在该圆轨道的$A$点第三次喷气，希望在对径方向的地表点$B$与地表相切。分别讨论向正前方排气和沿径向向外排气；检查两者是否都能到达指定的$B$。若只要求在某地表点相切而不限定方位，求两种比冲量，并在$r_0=2R_E$时比较。
\end{enumerate}
各次排气方向指相对飞船的喷气方向，与飞船获得的冲量方向相反。“与地表相切”指飞船到地表时的瞬时速度沿当地切线，落地速率由轨道确定。


\begin{center}\begin{minipage}{.98\linewidth}\centering
\begin{tikzpicture}[scale=.9,transform shape]
\draw[object] (0.00000,0.00000) circle(1.00000);\draw[guide] (0.00000,0.00000) circle(2.00000);\draw[guide] (0.00000,0.00000) circle(4.00000);\fill (0.00000,0.00000) circle(.035);\node[below] at (0.00000,0.00000) {$O$};\fill (4.00000,0.00000) circle(.035);\node[right] at (4.00000,0.00000) {$A$};\fill (-1.00000,0.00000) circle(.035);\node[left] at (-1.00000,0.00000) {$B$};\draw[guide] (-1.00000,0.00000)--(4.00000,0.00000);\node[above] at (0.00000,1.00000) {$R_E$};\node[above left] at (-2.40000,0.40000) {$r_0$};\node[above] at (0.00000,4.00000) {$r_1$};\draw[motion] (-2.00000,0.00000)--(-2.00000,-1.30000);\node[left] at (-2.10000,-1.00000) {$v_0$};
\end{tikzpicture}
\figcaption{图力试2.2\quad 三条同心尺度及第三次喷气的指定对径地表点B。}
\end{minipage}\end{center}

\Needspace{6\baselineskip}\subsection*{解答}


图力试2.2取地心为$O$。令$\mu=GM_E$，轨道速度在当地径向$\hat r$和沿运动切向$\hat t$中分解。喷气阶段可改变机械能与角动量；两次喷气之间只受中心引力，这两个量才守恒。

圆轨道的向心条件为$m v_0^2/r_0=\mu m/r_0^2$，飞船质量约去后确定初始圆速。第一次喷气的目标速率再由单位质量总能量为零确定，两个速度条件对应不同轨道。
（1）初始圆轨道速度为$v_0=\sqrt{\mu/r_0}$。刚能逃逸要求喷气后单位质量机械能为零：
\[
\frac{v_{e0}^2}{2}-\frac{\mu}{r_0}=0,\qquad
v_{e0}=\sqrt{\frac{2\mu}{r_0}}.
\]
在给定速度大小$v_{e0}$时，与初始速度同向的速度改变量最小，由速度矢量三角不等式得
\[
\boxed{\Delta v_0=v_{e0}-v_0
=(\sqrt2-1)\sqrt{\frac{\mu}{r_0}}}.
\]
飞船加速方向沿$\hat t$，所以气体向后喷出。随后飞船沿近地点半径$r_0$的抛物线向外运动。

\begin{center}\begin{minipage}{.98\linewidth}\centering
\begin{tikzpicture}[scale=.82,transform shape]
\draw[axis] (0.00000,0.00000)--(4.45000,0.00000);\draw[axis] (0.00000,0.00000)--(0.00000,3.65000);\node[below right] at (4.45000,0.00000) {$\hat r$};\node[above right] at (0.00000,3.65000) {$\hat t$};\draw[motion] (0.00000,0.00000)--(2.82843,2.00000);\node[above right] at (2.82843,2.00000) {$\boldsymbol v_e$};\draw[force] (0.00000,0.00000)--(0.00000,2.44949);\node[left] at (-0.10000,1.50000) {$v_c$};\draw[force] (2.82843,2.00000)--(0.00000,2.44949);\node[above] at (1.20000,2.80000) {$\Delta\boldsymbol v_1$};\draw[guide] (2.82843,2.00000)--(2.82843,0.00000);\node[below] at (1.41421,-0.10000) {$v_r$};\node[right] at (2.92843,1.00000) {$v_t$};\draw[motion] (5.00000,2.00000)--(6.60000,1.75000);\node[above] at (5.80000,2.10000) {排气 $-\Delta\boldsymbol v_1$};\node[above] at (3.20000,4.20000) {示例 $r_1=3r_0>2r_0$};\node[below] at (3.00000,-1.20000) {先消径向分量，再调切向分量};
\end{tikzpicture}
\figcaption{解图力试2.2A\quad 第二冲是速度矢量差，排气方向与它相反。}
\end{minipage}\end{center}
（2）解图力试2.2A在第二次喷气点画出此前的径向和切向速度。抛物线单位质量能量为零，角动量$j=\sqrt{2\mu r_0}$。在$r_1$处，
\[
v_e^2=\frac{2\mu}{r_1},\qquad
v_t=\frac{j}{r_1}=\frac{\sqrt{2\mu r_0}}{r_1},
\]
\[
v_r=\sqrt{v_e^2-v_t^2}
=\frac{\sqrt{2\mu(r_1-r_0)}}{r_1}>0.
\]
目标圆轨道速度为$v_c=\sqrt{\mu/r_1}$，只有切向分量。因此飞船所需比冲量矢量为
\[
\boxed{\Delta\boldsymbol v_1=-v_r\hat r+(v_c-v_t)\hat t},
\]
\[
\boxed{\Delta v_1
=\sqrt{\frac{\mu}{r_1}
\left(3-2\sqrt{\frac{2r_0}{r_1}}\right)}}.
\]
排气沿反方向$v_r\hat r+(v_t-v_c)\hat t$。其径向分量始终向外，而切向分量在$r_1=2r_0$变号：$r_1<2r_0$时沿前方，$r_1>2r_0$时沿后方，等号时纯径向向外。径向分量和切向分量一起确定入圆冲量与排气方向。

若将喷气前速度与圆轨道切向的夹角记为$\phi$，则$\cos\phi=v_t/v_e=\sqrt{r_0/r_1}$。速度三角形也给$\Delta v_1^2=v_e^2+v_c^2-2v_ev_c\cos\phi$，其中径向部分已由完整夹角计入。比较排气方向时，下面的点积式又可写成余弦定理形式$\cos\alpha=(v_e^2+\Delta v_1^2-v_c^2)/(2v_e\Delta v_1)$，两种表达使用同一组速度矢量。
若以$\alpha$表示排气与喷气前飞船速度的夹角，由点积可得
\[
\cos\alpha=
\frac{v_r^2+(v_t-v_c)v_t}{v_e\Delta v_1}
=\frac{2-\sqrt{2r_0/r_1}}
{\sqrt{2(3-2\sqrt{2r_0/r_1})}}.
\]

\begin{center}\begin{minipage}{.98\linewidth}\centering
\begin{tikzpicture}[scale=.83,transform shape]
\draw[object] (0.00000,0.00000) circle(1.00000);\node[anchor=center] at (0.00000,0.00000) {$O$};\draw[traj] plot coordinates {(4.00000,0.00000) (3.99897,0.05737) (3.99589,0.11470) (3.99075,0.17193) (3.98356,0.22902) (3.97432,0.28592) (3.96305,0.34259) (3.94975,0.39897) (3.93444,0.45503) (3.91712,0.51071) (3.89781,0.56597) (3.87653,0.62076) (3.85329,0.67505) (3.82812,0.72877) (3.80103,0.78190) (3.77204,0.83439) (3.74119,0.88619) (3.70849,0.93725) (3.67397,0.98755) (3.63767,1.03704) (3.59960,1.08567) (3.55981,1.13340) (3.51832,1.18021) (3.47517,1.22604) (3.43039,1.27086) (3.38403,1.31464) (3.33611,1.35734) (3.28668,1.39892) (3.23579,1.43934) (3.18346,1.47859) (3.12975,1.51661) (3.07469,1.55339) (3.01834,1.58889) (2.96075,1.62308) (2.90194,1.65593) (2.84199,1.68742) (2.78093,1.71753) (2.71881,1.74622) (2.65570,1.77347) (2.59163,1.79926) (2.52666,1.82357) (2.46085,1.84638) (2.39425,1.86767) (2.32691,1.88743) (2.25889,1.90563) (2.19024,1.92226) (2.12103,1.93731) (2.05131,1.95076) (1.98113,1.96261) (1.91056,1.97285) (1.83964,1.98146) (1.76845,1.98844) (1.69704,1.99378) (1.62547,1.99748) (1.55379,1.99954) (1.48207,1.99995) (1.41036,1.99871) (1.33873,1.99583) (1.26723,1.99131) (1.19592,1.98515) (1.12486,1.97736) (1.05411,1.96793) (0.98373,1.95689) (0.91377,1.94424) (0.84430,1.92998) (0.77536,1.91414) (0.70702,1.89672) (0.63934,1.87774) (0.57236,1.85722) (0.50614,1.83517) (0.44075,1.81160) (0.37622,1.78655) (0.31262,1.76002) (0.25000,1.73205) (0.18841,1.70265) (0.12789,1.67185) (0.06851,1.63967) (0.01030,1.60615) (-0.04668,1.57130) (-0.10239,1.53516) (-0.15677,1.49775) (-0.20980,1.45912) (-0.26142,1.41928) (-0.31158,1.37827) (-0.36026,1.33613) (-0.40741,1.29289) (-0.45298,1.24858) (-0.49695,1.20325) (-0.53927,1.15692) (-0.57992,1.10965) (-0.61885,1.06146) (-0.65604,1.01240) (-0.69146,0.96250) (-0.72507,0.91181) (-0.75685,0.86037) (-0.78677,0.80823) (-0.81481,0.75542) (-0.84095,0.70198) (-0.86516,0.64797) (-0.88742,0.59343) (-0.90771,0.53839) (-0.92603,0.48292) (-0.94235,0.42704) (-0.95665,0.37082) (-0.96894,0.31429) (-0.97919,0.25750) (-0.98741,0.20050) (-0.99357,0.14333) (-0.99769,0.08604) (-0.99974,0.02869) (-0.99974,-0.02869) (-0.99769,-0.08604) (-0.99357,-0.14333) (-0.98741,-0.20050) (-0.97919,-0.25750) (-0.96894,-0.31429) (-0.95665,-0.37082) (-0.94235,-0.42704) (-0.92603,-0.48292) (-0.90771,-0.53839) (-0.88742,-0.59343) (-0.86516,-0.64797) (-0.84095,-0.70198) (-0.81481,-0.75542) (-0.78677,-0.80823) (-0.75685,-0.86037) (-0.72507,-0.91181) (-0.69146,-0.96250) (-0.65604,-1.01240) (-0.61885,-1.06146) (-0.57992,-1.10965) (-0.53927,-1.15692) (-0.49695,-1.20325) (-0.45298,-1.24858) (-0.40741,-1.29289) (-0.36026,-1.33613) (-0.31158,-1.37827) (-0.26142,-1.41928) (-0.20980,-1.45912) (-0.15677,-1.49775) (-0.10239,-1.53516) (-0.04668,-1.57130) (0.01030,-1.60615) (0.06851,-1.63967) (0.12789,-1.67185) (0.18841,-1.70265) (0.25000,-1.73205) (0.31262,-1.76002) (0.37622,-1.78655) (0.44075,-1.81160) (0.50614,-1.83517) (0.57236,-1.85722) (0.63934,-1.87774) (0.70702,-1.89672) (0.77536,-1.91414) (0.84430,-1.92998) (0.91377,-1.94424) (0.98373,-1.95689) (1.05411,-1.96793) (1.12486,-1.97736) (1.19592,-1.98515) (1.26723,-1.99131) (1.33873,-1.99583) (1.41036,-1.99871) (1.48207,-1.99995) (1.55379,-1.99954) (1.62547,-1.99748) (1.69704,-1.99378) (1.76845,-1.98844) (1.83964,-1.98146) (1.91056,-1.97285) (1.98113,-1.96261) (2.05131,-1.95076) (2.12103,-1.93731) (2.19024,-1.92226) (2.25889,-1.90563) (2.32691,-1.88743) (2.39425,-1.86767) (2.46085,-1.84638) (2.52666,-1.82357) (2.59163,-1.79926) (2.65570,-1.77347) (2.71881,-1.74622) (2.78093,-1.71753) (2.84199,-1.68742) (2.90194,-1.65593) (2.96075,-1.62308) (3.01834,-1.58889) (3.07469,-1.55339) (3.12975,-1.51661) (3.18346,-1.47859) (3.23579,-1.43934) (3.28668,-1.39892) (3.33611,-1.35734) (3.38403,-1.31464) (3.43039,-1.27086) (3.47517,-1.22604) (3.51832,-1.18021) (3.55981,-1.13340) (3.59960,-1.08567) (3.63767,-1.03704) (3.67397,-0.98755) (3.70849,-0.93725) (3.74119,-0.88619) (3.77204,-0.83439) (3.80103,-0.78190) (3.82812,-0.72877) (3.85329,-0.67505) (3.87653,-0.62076) (3.89781,-0.56597) (3.91712,-0.51071) (3.93444,-0.45503) (3.94975,-0.39897) (3.96305,-0.34259) (3.97432,-0.28592) (3.98356,-0.22902) (3.99075,-0.17193) (3.99589,-0.11470) (3.99897,-0.05737) (4.00000,-0.00000)};\fill (4.00000,0.00000) circle(.035);\node[below right] at (4.00000,0.00000) {$A$};\fill (-1.00000,0.00000) circle(.035);\node[left] at (-1.00000,0.00000) {$B$};\draw[guide] (-1.00000,0.00000)--(4.00000,0.00000);\draw[motion] (4.00000,0.00000)--(4.00000,1.10000);\node[above left] at (4.00000,1.10000) {$v_a$};\draw[force] (4.60000,0.00000)--(4.60000,1.70000);\node[right] at (4.80000,1.70000) {前向排气};\draw[motion] (-1.00000,0.00000)--(-1.00000,-1.50000);\node[below left] at (-1.00000,-1.50000) {$v_p$};\node[above] at (1.50000,3.00000) {A远地点，B近地点，方位差 $\pi$};\draw[<->,thin,color=muted] (0.00000,-3.00000)--(4.00000,-3.00000);\node[below] at (2.00000,-3.00000) {$r_1$};\draw[<->,thin,color=muted] (-1.00000,-4.00000)--(0.00000,-4.00000);\node[below] at (-0.50000,-4.00000) {$R_E$};
\end{tikzpicture}
\figcaption{解图力试2.2B\quad 切向减速产生对径拱点，能到达指定B。}
\end{minipage}\end{center}
（3）解图力试2.2B的切向方案向前排气，飞船切向减速。欲从$A$到对径地表点$B$并相切，$A$为椭圆远地点、$B$为近地点。设减速后速率为$v_a$、到$B$时为$v_p$，则
\[
r_1v_a=R_Ev_p,\qquad
\frac{v_a^2}{2}-\frac{\mu}{r_1}
=\frac{v_p^2}{2}-\frac{\mu}{R_E}.
\]
由角动量式先取$v_p=(r_1/R_E)v_a$，代入能量式有$v_a^2[(r_1/R_E)^2-1]=2\mu(1/R_E-1/r_1)$。两边含有的$r_1-R_E$因子约去后得到远地点速率，取正根对应仍沿原方向运动的返回支。
消去$v_p$，
\[
v_a^2=\frac{2\mu R_E}{r_1(r_1+R_E)},\qquad
\boxed{\Delta v_t=v_c-v_a
=v_c\left(1-\sqrt{\frac{2R_E}{r_1+R_E}}\right)}.
\]
该椭圆的两拱点确在同一直径两侧，符合指定落地点。

\begin{center}\begin{minipage}{.98\linewidth}\centering
\begin{tikzpicture}[scale=.83,transform shape]
\draw[object] (0.00000,0.00000) circle(1.00000);\node[below right] at (0.00000,-0.30000) {$O$};\draw[traj] plot coordinates {(3.00000,0.00000) (2.92290,0.03842) (2.84917,0.07492) (2.77858,0.10963) (2.71091,0.14267) (2.64597,0.17415) (2.58357,0.20419) (2.52356,0.23286) (2.46579,0.26026) (2.41011,0.28646) (2.35641,0.31154) (2.30457,0.33556) (2.25447,0.35859) (2.20603,0.38068) (2.15914,0.40189) (2.11373,0.42225) (2.06971,0.44183) (2.02701,0.46065) (1.98557,0.47876) (1.94531,0.49620) (1.90618,0.51300) (1.86812,0.52919) (1.83109,0.54479) (1.79502,0.55985) (1.75988,0.57438) (1.72563,0.58841) (1.69221,0.60196) (1.65960,0.61505) (1.62776,0.62770) (1.59665,0.63994) (1.56624,0.65178) (1.53651,0.66323) (1.50742,0.67432) (1.47895,0.68505) (1.45107,0.69544) (1.42376,0.70551) (1.39699,0.71527) (1.37074,0.72473) (1.34500,0.73390) (1.31973,0.74279) (1.29494,0.75142) (1.27058,0.75979) (1.24666,0.76791) (1.22315,0.77579) (1.20004,0.78343) (1.17731,0.79086) (1.15495,0.79807) (1.13295,0.80507) (1.11129,0.81187) (1.08996,0.81847) (1.06895,0.82489) (1.04824,0.83112) (1.02783,0.83718) (1.00771,0.84307) (0.98787,0.84879) (0.96829,0.85434) (0.94897,0.85974) (0.92990,0.86499) (0.91107,0.87009) (0.89248,0.87505) (0.87411,0.87987) (0.85595,0.88455) (0.83801,0.88911) (0.82027,0.89353) (0.80272,0.89783) (0.78537,0.90200) (0.76820,0.90606) (0.75121,0.91000) (0.73438,0.91383) (0.71773,0.91754) (0.70123,0.92115) (0.68489,0.92466) (0.66870,0.92806) (0.65265,0.93136) (0.63675,0.93457) (0.62098,0.93767) (0.60534,0.94069) (0.58982,0.94361) (0.57443,0.94644) (0.55916,0.94918) (0.54400,0.95184) (0.52895,0.95441) (0.51400,0.95690) (0.49916,0.95930) (0.48441,0.96163) (0.46976,0.96387) (0.45521,0.96604) (0.44073,0.96813) (0.42635,0.97015) (0.41204,0.97209) (0.39782,0.97396) (0.38366,0.97576) (0.36958,0.97749) (0.35557,0.97915) (0.34162,0.98073) (0.32774,0.98226) (0.31392,0.98371) (0.30015,0.98510) (0.28644,0.98642) (0.27278,0.98767) (0.25917,0.98887) (0.24561,0.99000) (0.23209,0.99106) (0.21861,0.99207) (0.20517,0.99301) (0.19176,0.99389) (0.17839,0.99471) (0.16505,0.99547) (0.15174,0.99617) (0.13846,0.99681) (0.12520,0.99739) (0.11196,0.99791) (0.09875,0.99838) (0.08554,0.99878) (0.07236,0.99913) (0.05919,0.99942) (0.04602,0.99965) (0.03287,0.99982) (0.01972,0.99994) (0.00657,0.99999) (-0.00657,0.99999) (-0.01972,0.99994) (-0.03287,0.99982) (-0.04602,0.99965) (-0.05919,0.99942) (-0.07236,0.99913) (-0.08554,0.99878) (-0.09875,0.99838) (-0.11196,0.99791) (-0.12520,0.99739) (-0.13846,0.99681) (-0.15174,0.99617) (-0.16505,0.99547) (-0.17839,0.99471) (-0.19176,0.99389) (-0.20517,0.99301) (-0.21861,0.99207) (-0.23209,0.99106) (-0.24561,0.99000) (-0.25917,0.98887) (-0.27278,0.98767) (-0.28644,0.98642) (-0.30015,0.98510) (-0.31392,0.98371) (-0.32774,0.98226) (-0.34162,0.98073) (-0.35557,0.97915) (-0.36958,0.97749) (-0.38366,0.97576) (-0.39782,0.97396) (-0.41204,0.97209) (-0.42635,0.97015) (-0.44073,0.96813) (-0.45521,0.96604) (-0.46976,0.96387) (-0.48441,0.96163) (-0.49916,0.95930) (-0.51400,0.95690) (-0.52895,0.95441) (-0.54400,0.95184) (-0.55916,0.94918) (-0.57443,0.94644) (-0.58982,0.94361) (-0.60534,0.94069) (-0.62098,0.93767) (-0.63675,0.93457) (-0.65265,0.93136) (-0.66870,0.92806) (-0.68489,0.92466) (-0.70123,0.92115) (-0.71773,0.91754) (-0.73438,0.91383) (-0.75121,0.91000) (-0.76820,0.90606) (-0.78537,0.90200) (-0.80272,0.89783) (-0.82027,0.89353) (-0.83801,0.88911) (-0.85595,0.88455) (-0.87411,0.87987) (-0.89248,0.87505) (-0.91107,0.87009) (-0.92990,0.86499) (-0.94897,0.85974) (-0.96829,0.85434) (-0.98787,0.84879) (-1.00771,0.84307) (-1.02783,0.83718) (-1.04824,0.83112) (-1.06895,0.82489) (-1.08996,0.81847) (-1.11129,0.81187) (-1.13295,0.80507) (-1.15495,0.79807) (-1.17731,0.79086) (-1.20004,0.78343) (-1.22315,0.77579) (-1.24666,0.76791) (-1.27058,0.75979) (-1.29494,0.75142) (-1.31973,0.74279) (-1.34500,0.73390) (-1.37074,0.72473) (-1.39699,0.71527) (-1.42376,0.70551) (-1.45107,0.69544) (-1.47895,0.68505) (-1.50742,0.67432) (-1.53651,0.66323) (-1.56624,0.65178) (-1.59665,0.63994) (-1.62776,0.62770) (-1.65960,0.61505) (-1.69221,0.60196) (-1.72563,0.58841) (-1.75988,0.57438) (-1.79502,0.55985) (-1.83109,0.54479) (-1.86812,0.52919) (-1.90618,0.51300) (-1.94531,0.49620) (-1.98557,0.47876) (-2.02701,0.46065) (-2.06971,0.44183) (-2.11373,0.42225) (-2.15914,0.40189) (-2.20603,0.38068) (-2.25447,0.35859) (-2.30457,0.33556) (-2.35641,0.31154) (-2.41011,0.28646) (-2.46579,0.26026) (-2.52356,0.23286) (-2.58357,0.20419) (-2.64597,0.17415) (-2.71091,0.14267) (-2.77858,0.10963) (-2.84917,0.07492) (-2.92290,0.03842) (-3.00000,0.00000)};\fill (3.00000,0.00000) circle(.035);\node[below right] at (3.00000,0.00000) {$A$};\fill (0.00000,1.00000) circle(.035);\node[above] at (0.00000,1.00000) {$B_r$相切};\fill (-1.00000,0.00000) circle(.035);\node[below left] at (-1.00000,0.00000) {指定 $B$};\fill (-3.00000,0.00000) circle(.035);\node[below left] at (-3.00000,0.00000) {$r(\pi)=r_1$};\draw[guide] (-3.40000,0.00000)--(4.80000,0.00000);\draw[force] (3.00000,0.00000)--(1.40000,0.80000);\node[above right] at (1.40000,0.80000) {$-u\hat r+v_c\hat t$};\draw[motion] (3.00000,-0.60000)--(4.60000,-0.60000);\node[below] at (4.00000,-0.70000) {径向向外排气};\node[above] at (0.00000,2.60000) {示例 $r_1=3R_E$：双曲线先向内运行};\node[below] at (0.00000,-2.30000) {相切点在90°方位，对径方位处半径仍为$r_1$};
\end{tikzpicture}
\figcaption{解图力试2.2C\quad 纯径向冲量的近地点在Br，与指定B方位不同。}
\end{minipage}\end{center}
解图力试2.2C的径向排气给飞船一个向内比冲量$u$，切向速度仍为$v_c$，单位质量角动量仍为$j=r_1v_c$。若暂时只要求在某地表点$B_r$相切，能量与角动量给
\[
\frac{v_c^2+u^2}{2}-\frac{\mu}{r_1}
=\frac{(r_1v_c/R_E)^2}{2}-\frac{\mu}{R_E},
\]
\[
\boxed{u=\Delta v_r=v_c\left(\frac{r_1}{R_E}-1\right)}.
\]
接着由同一喷后速度求近地点方位，将接地半径与指定位置一起确定。把$A$放在$+x$轴，取初始切向速度向$+y$，径向喷气后的速度是$-u\hat x+v_c\hat y$。这里$\boldsymbol j=r_1v_c\hat z$，因此$\boldsymbol v\times\boldsymbol j=r_1v_c^2\hat x+r_1uv_c\hat y$。第一项除以$\mu$恰为$\hat x$，与偏心率定义中的径向单位矢量抵消，余下分量确定真正的近地点方向。
它的轨道偏心率矢量为
\[
\boldsymbol e=\frac{\boldsymbol v\times\boldsymbol j}{\mu}-\hat r
=\frac{u}{v_c}\hat y.
\]
近地点沿$+y$，离$A$的方位差为$\pi/2$，不是对径的$\pi$。轨道方程也直接说明：
\[
r(\phi)=\frac{r_1}{1+(u/v_c)\sin\phi},
\qquad r(\pi)=r_1>R_E.
\]
所以\textbf{纯径向方案的近地点在$B_r$，对径方位处的轨道半径为$r_1>R_E$}。上面的$\Delta v_r$只能回答放宽落地点后的相切问题；该轨道满足接地半径，近地点方位则由偏心率矢量指向$B_r$。

对于放宽的任务，令$q=r_1/R_E>1$，则轨道偏心率$e=q-1$。当$1<q<2$时为椭圆，$q=2$时为抛物线，$q>2$时为先向内接近地球的双曲线。若$r_0=2R_E$且$r_1>r_0$，径向返回段实际属于双曲线。

最后比较两种放宽方案的比冲量：
\[
\frac{\Delta v_r-\Delta v_t}{v_c}
=q-2+\sqrt{\frac2{q+1}}>0,\qquad q>1.
\]
右式在$q=1$为零，导数为$1-\tfrac1{\sqrt2}(q+1)^{-3/2}>0$，因此
\[
\boxed{\Delta v_r>\Delta v_t}.
\]
这个比较对$r_0=2R_E$成立，实际上对任意$r_1>R_E$都成立；但二者相切的地表点不同，比较的是两条在不同地表位置相切的轨道。
\begin{center}\begin{minipage}{.98\linewidth}\centering
\begin{tikzpicture}[scale=.8,transform shape]
\draw[axis] (0,0)--(10.7,0) node[right] {$q=r_1/R_E$};
\draw[axis] (0,0)--(0,3.2) node[above] {第三冲大小$/v_c$};
\node[above] at (5,3.5) {放宽地表相切点方位后：实线切向方案，虚线径向方案};
\draw[traj] plot coordinates {(0.003333333,0.000224916) (0.086638889,0.005791734) (0.169944444,0.011256491) (0.253250000,0.016622268) (0.336555556,0.021892015) (0.419861111,0.027068565) (0.503166667,0.032154632) (0.586472222,0.037152822) (0.669777778,0.042065636) (0.753083333,0.046895479) (0.836388889,0.051644658) (0.919694444,0.056315395) (1.003000000,0.060909826) (1.086305556,0.065430006) (1.169611111,0.069877914) (1.252916667,0.074255455) (1.336222222,0.078564465) (1.419527778,0.082806714) (1.502833333,0.086983909) (1.586138889,0.091097694) (1.669444444,0.095149659) (1.752750000,0.099141336) (1.836055556,0.103074205) (1.919361111,0.106949697) (2.002666667,0.110769193) (2.085972222,0.114534030) (2.169277778,0.118245497) (2.252583333,0.121904844) (2.335888889,0.125513281) (2.419194444,0.129071976) (2.502500000,0.132582062) (2.585805556,0.136044635) (2.669111111,0.139460759) (2.752416667,0.142831461) (2.835722222,0.146157740) (2.919027778,0.149440564) (3.002333333,0.152680870) (3.085638889,0.155879567) (3.168944444,0.159037540) (3.252250000,0.162155645) (3.335555556,0.165234713) (3.418861111,0.168275553) (3.502166667,0.171278949) (3.585472222,0.174245664) (3.668777778,0.177176437) (3.752083333,0.180071989) (3.835388889,0.182933020) (3.918694444,0.185760210) (4.002000000,0.188554221) (4.085305556,0.191315698) (4.168611111,0.194045266) (4.251916667,0.196743535) (4.335222222,0.199411099) (4.418527778,0.202048537) (4.501833333,0.204656410) (4.585138889,0.207235267) (4.668444444,0.209785642) (4.751750000,0.212308056) (4.835055556,0.214803016) (4.918361111,0.217271016) (5.001666667,0.219712539) (5.084972222,0.222128054) (5.168277778,0.224518019) (5.251583333,0.226882883) (5.334888889,0.229223082) (5.418194444,0.231539040) (5.501500000,0.233831175) (5.584805556,0.236099891) (5.668111111,0.238345585) (5.751416667,0.240568643) (5.834722222,0.242769443) (5.918027778,0.244948354) (6.001333333,0.247105737) (6.084638889,0.249241943) (6.167944444,0.251357318) (6.251250000,0.253452197) (6.334555556,0.255526908) (6.417861111,0.257581775) (6.501166667,0.259617110) (6.584472222,0.261633223) (6.667777778,0.263630412) (6.751083333,0.265608972) (6.834388889,0.267569192) (6.917694444,0.269511353) (7.001000000,0.271435730) (7.084305556,0.273342593) (7.167611111,0.275232206) (7.250916667,0.277104828) (7.334222222,0.278960712) (7.417527778,0.280800105) (7.500833333,0.282623251) (7.584138889,0.284430388) (7.667444444,0.286221747) (7.750750000,0.287997558) (7.834055556,0.289758044) (7.917361111,0.291503424) (8.000666667,0.293233914) (8.083972222,0.294949723) (8.167277778,0.296651058) (8.250583333,0.298338121) (8.333888889,0.300011111) (8.417194444,0.301670222) (8.500500000,0.303315646) (8.583805556,0.304947569) (8.667111111,0.306566174) (8.750416667,0.308171643) (8.833722222,0.309764152) (8.917027778,0.311343874) (9.000333333,0.312910980) (9.083638889,0.314465636) (9.166944444,0.316008007) (9.250250000,0.317538253) (9.333555556,0.319056532) (9.416861111,0.320563001) (9.500166667,0.322057810) (9.583472222,0.323541110) (9.666777778,0.325013047) (9.750083333,0.326473767) (9.833388889,0.327923410) (9.916694444,0.329362115) (10.000000000,0.330790021)};
\draw[dashed,thick,color=warm] plot coordinates {(0.003333333,0.000900000) (0.086638889,0.023392500) (0.169944444,0.045885000) (0.253250000,0.068377500) (0.336555556,0.090870000) (0.419861111,0.113362500) (0.503166667,0.135855000) (0.586472222,0.158347500) (0.669777778,0.180840000) (0.753083333,0.203332500) (0.836388889,0.225825000) (0.919694444,0.248317500) (1.003000000,0.270810000) (1.086305556,0.293302500) (1.169611111,0.315795000) (1.252916667,0.338287500) (1.336222222,0.360780000) (1.419527778,0.383272500) (1.502833333,0.405765000) (1.586138889,0.428257500) (1.669444444,0.450750000) (1.752750000,0.473242500) (1.836055556,0.495735000) (1.919361111,0.518227500) (2.002666667,0.540720000) (2.085972222,0.563212500) (2.169277778,0.585705000) (2.252583333,0.608197500) (2.335888889,0.630690000) (2.419194444,0.653182500) (2.502500000,0.675675000) (2.585805556,0.698167500) (2.669111111,0.720660000) (2.752416667,0.743152500) (2.835722222,0.765645000) (2.919027778,0.788137500) (3.002333333,0.810630000) (3.085638889,0.833122500) (3.168944444,0.855615000) (3.252250000,0.878107500) (3.335555556,0.900600000) (3.418861111,0.923092500) (3.502166667,0.945585000) (3.585472222,0.968077500) (3.668777778,0.990570000) (3.752083333,1.013062500) (3.835388889,1.035555000) (3.918694444,1.058047500) (4.002000000,1.080540000) (4.085305556,1.103032500) (4.168611111,1.125525000) (4.251916667,1.148017500) (4.335222222,1.170510000) (4.418527778,1.193002500) (4.501833333,1.215495000) (4.585138889,1.237987500) (4.668444444,1.260480000) (4.751750000,1.282972500) (4.835055556,1.305465000) (4.918361111,1.327957500) (5.001666667,1.350450000) (5.084972222,1.372942500) (5.168277778,1.395435000) (5.251583333,1.417927500) (5.334888889,1.440420000) (5.418194444,1.462912500) (5.501500000,1.485405000) (5.584805556,1.507897500) (5.668111111,1.530390000) (5.751416667,1.552882500) (5.834722222,1.575375000) (5.918027778,1.597867500) (6.001333333,1.620360000) (6.084638889,1.642852500) (6.167944444,1.665345000) (6.251250000,1.687837500) (6.334555556,1.710330000) (6.417861111,1.732822500) (6.501166667,1.755315000) (6.584472222,1.777807500) (6.667777778,1.800300000) (6.751083333,1.822792500) (6.834388889,1.845285000) (6.917694444,1.867777500) (7.001000000,1.890270000) (7.084305556,1.912762500) (7.167611111,1.935255000) (7.250916667,1.957747500) (7.334222222,1.980240000) (7.417527778,2.002732500) (7.500833333,2.025225000) (7.584138889,2.047717500) (7.667444444,2.070210000) (7.750750000,2.092702500) (7.834055556,2.115195000) (7.917361111,2.137687500) (8.000666667,2.160180000) (8.083972222,2.182672500) (8.167277778,2.205165000) (8.250583333,2.227657500) (8.333888889,2.250150000) (8.417194444,2.272642500) (8.500500000,2.295135000) (8.583805556,2.317627500) (8.667111111,2.340120000) (8.750416667,2.362612500) (8.833722222,2.385105000) (8.917027778,2.407597500) (9.000333333,2.430090000) (9.083638889,2.452582500) (9.166944444,2.475075000) (9.250250000,2.497567500) (9.333555556,2.520060000) (9.416861111,2.542552500) (9.500166667,2.565045000) (9.583472222,2.587537500) (9.666777778,2.610030000) (9.750083333,2.632522500) (9.833388889,2.655015000) (9.916694444,2.677507500) (10.000000000,2.700000000)};
\node[below] at (0,-.12) {$1$};
\node[below] at (3.3333333333333335,-.12) {$2$};
\node[below] at (10,-.12) {$4$};
\end{tikzpicture}
\figcaption{解图力试2.2D\quad 放宽地表相切点方位后：实线切向方案，虚线径向方案}
\end{minipage}\end{center}
图D的两条曲线分别对应前述对径椭圆和在$B_r$相切的径向返回轨道。它们可以比较冲量大小，同时各自保持由偏心率矢量确定的空间方位。


\Needspace{13\baselineskip}

\pdfbookmark[2]{第 3 题 绳越过凸圆弧时的张力与脱离}{q3}

\section*{第 3 题\quad 绳越过凸圆弧时的张力与脱离}\label{q:3}


一条均匀柔绳总长为$l$、线密度为$\lambda$，放在光滑水平面上，右端位于与下降斜面相接的凸边$B$。斜面倾角为$\theta$，$0<\theta<\pi/2$；连接处用半径$r\ll l$的小圆弧平滑过渡。绳可自由弯曲且不可伸长。

让绳端因微小扰动先下滑一小段，随后在重力下运动。以$x$表示尚在水平面上的绳长。求下滑过程中的绳张力、凸边对绳的支持力，以及绳开始脱离凸边时的位置。先求$r/l\to0$的领先阶结果，并说明有限$r$时脱离条件如何修正。圆弧没有黏附能力，只能向外推绳。


\begin{center}\begin{minipage}{.98\linewidth}\centering
\begin{tikzpicture}[scale=.9,transform shape]
\draw[object] (-4.00000,0.00000)--(0.00000,0.00000);\draw[object] plot coordinates {(0.00000,-0.00000) (0.00740,-0.00003) (0.01481,-0.00012) (0.02221,-0.00027) (0.02961,-0.00049) (0.03701,-0.00076) (0.04441,-0.00110) (0.05181,-0.00149) (0.05920,-0.00195) (0.06658,-0.00247) (0.07397,-0.00304) (0.08134,-0.00368) (0.08872,-0.00438) (0.09608,-0.00514) (0.10344,-0.00596) (0.11079,-0.00685) (0.11814,-0.00779) (0.12548,-0.00879) (0.13280,-0.00985) (0.14012,-0.01098) (0.14743,-0.01216) (0.15473,-0.01340) (0.16202,-0.01470) (0.16930,-0.01607) (0.17657,-0.01749) (0.18382,-0.01897) (0.19107,-0.02052) (0.19830,-0.02212) (0.20551,-0.02378) (0.21271,-0.02550) (0.21990,-0.02728) (0.22708,-0.02912) (0.23423,-0.03102) (0.24138,-0.03297) (0.24850,-0.03499) (0.25561,-0.03706) (0.26270,-0.03919) (0.26977,-0.04138) (0.27683,-0.04363) (0.28387,-0.04594) (0.29088,-0.04830) (0.29788,-0.05073) (0.30486,-0.05321) (0.31182,-0.05574) (0.31875,-0.05834) (0.32567,-0.06099) (0.33256,-0.06370) (0.33943,-0.06646) (0.34627,-0.06928) (0.35310,-0.07216) (0.35990,-0.07509) (0.36667,-0.07808) (0.37342,-0.08113) (0.38015,-0.08423) (0.38685,-0.08738) (0.39352,-0.09059) (0.40017,-0.09386) (0.40678,-0.09718) (0.41338,-0.10055) (0.41994,-0.10398) (0.42648,-0.10746) (0.43298,-0.11100) (0.43946,-0.11459) (0.44591,-0.11823) (0.45232,-0.12192) (0.45871,-0.12567) (0.46507,-0.12947) (0.47139,-0.13332) (0.47768,-0.13723) (0.48394,-0.14118) (0.49017,-0.14519) (0.49636,-0.14925) (0.50252,-0.15336) (0.50865,-0.15752) (0.51474,-0.16173) (0.52080,-0.16599) (0.52682,-0.17030) (0.53280,-0.17466) (0.53875,-0.17907) (0.54467,-0.18352)};\draw[object] (0.54467,-0.18352)--(3.72900,-2.60427);\draw[very thick,color=accent] (-3.00000,0.00000)--(0.00000,0.00000);\draw[traj] plot coordinates {(0.00000,-0.00000) (0.00740,-0.00003) (0.01481,-0.00012) (0.02221,-0.00027) (0.02961,-0.00049) (0.03701,-0.00076) (0.04441,-0.00110) (0.05181,-0.00149) (0.05920,-0.00195) (0.06658,-0.00247) (0.07397,-0.00304) (0.08134,-0.00368) (0.08872,-0.00438) (0.09608,-0.00514) (0.10344,-0.00596) (0.11079,-0.00685) (0.11814,-0.00779) (0.12548,-0.00879) (0.13280,-0.00985) (0.14012,-0.01098) (0.14743,-0.01216) (0.15473,-0.01340) (0.16202,-0.01470) (0.16930,-0.01607) (0.17657,-0.01749) (0.18382,-0.01897) (0.19107,-0.02052) (0.19830,-0.02212) (0.20551,-0.02378) (0.21271,-0.02550) (0.21990,-0.02728) (0.22708,-0.02912) (0.23423,-0.03102) (0.24138,-0.03297) (0.24850,-0.03499) (0.25561,-0.03706) (0.26270,-0.03919) (0.26977,-0.04138) (0.27683,-0.04363) (0.28387,-0.04594) (0.29088,-0.04830) (0.29788,-0.05073) (0.30486,-0.05321) (0.31182,-0.05574) (0.31875,-0.05834) (0.32567,-0.06099) (0.33256,-0.06370) (0.33943,-0.06646) (0.34627,-0.06928) (0.35310,-0.07216) (0.35990,-0.07509) (0.36667,-0.07808) (0.37342,-0.08113) (0.38015,-0.08423) (0.38685,-0.08738) (0.39352,-0.09059) (0.40017,-0.09386) (0.40678,-0.09718) (0.41338,-0.10055) (0.41994,-0.10398) (0.42648,-0.10746) (0.43298,-0.11100) (0.43946,-0.11459) (0.44591,-0.11823) (0.45232,-0.12192) (0.45871,-0.12567) (0.46507,-0.12947) (0.47139,-0.13332) (0.47768,-0.13723) (0.48394,-0.14118) (0.49017,-0.14519) (0.49636,-0.14925) (0.50252,-0.15336) (0.50865,-0.15752) (0.51474,-0.16173) (0.52080,-0.16599) (0.52682,-0.17030) (0.53280,-0.17466) (0.53875,-0.17907) (0.54467,-0.18352)};\draw[traj] (0.54467,-0.18352)--(2.77370,-1.87805);\fill (0.00000,0.00000) circle(.035);\node[above left] at (0.00000,0.00000) {$B$};\fill (0.00000,-0.90000) circle(.035);\node[below left] at (0.00000,-0.90000) {曲率中心};\draw[guide] (0.00000,-0.90000)--(0.54467,-0.18352);\node[right] at (0.20000,-0.55000) {$r$};\draw[<->,thin,color=muted] (-3.00000,1.00000)--(0.00000,1.00000);\node[below] at (-1.50000,1.00000) {$x$};\node[above] at (2.70000,-0.30000) {下降斜面 $\theta$};\node[above] at (0.00000,2.10000) {圆角放大；$r\ll l$};
\end{tikzpicture}
\figcaption{图力试2.3\quad 凸边半径很小，几何图将圆角放大显示。}
\end{minipage}\end{center}

\Needspace{6\baselineskip}\subsection*{解答}


图力试2.3中，绳沿水平段、凸圆弧、斜面段依次向右下运动。先在$r/l\to0$的领先阶忽略圆弧占据的长度与重力，保留其转向角$\theta$；水平绳长为$x$，斜面绳长为$l-x$。绳不可伸长，各段速率同为$v$，且$\dot x=-v$。

（1）取初始水平位置为势能零点。斜面绳质心低于水平面$(l-x)\sin\theta/2$，由能量守恒，
\[
\frac12\lambda l v^2
=\frac12\lambda g(l-x)^2\sin\theta,
\quad
\boxed{v^2=\frac{g\sin\theta}{l}(l-x)^2}.
\]
求时间导数，利用$\dot x=-v$，得
\[
a=\dot v=\frac{g\sin\theta}{l}(l-x).
\]
水平段的物质不断流出，流出质量率为$\lambda v$，流出动量率为$\lambda v^2$。因此水平段的动量方程为
\[
\frac{\dd}{\dd t}(\lambda xv)=T_1-\lambda v^2,
\quad
\lambda xa-\lambda v^2=T_1-\lambda v^2.
\]
流出项与质量变化项抵消，故$T_1=\lambda x a$。

\begin{center}\begin{minipage}{.98\linewidth}\centering
\begin{tikzpicture}[scale=.8,transform shape]
\draw[very thick,color=accent] (-3.00000,0.00000)--(0.00000,0.00000);\draw[force] (0.00000,0.00000)--(1.30000,0.00000);\node[above] at (0.80000,0.10000) {$T_1$};\draw[motion] (-2.00000,1.00000)--(-0.50000,1.00000);\node[above] at (-1.20000,1.10000) {$v$};\node[below] at (-1.50000,-0.30000) {水平段 $\lambda x$};\begin{scope}[shift={(7.00000,0.00000)}]
\draw[very thick,color=accent] (0.00000,0.00000)--(2.78629,-2.11815);\draw[force] (0.00000,0.00000)--(-0.95530,0.72622);\node[above left] at (-1.00000,0.70000) {$T_2$};\draw[force] (1.35334,-1.02882)--(1.35334,-2.42882);\node[below right] at (1.45334,-2.62882) {$\lambda(l-x)g$};\node[above] at (1.80000,1.40000) {斜段 $\lambda(l-x)$};
\end{scope}\node[below] at (4.00000,-4.00000) {两开口处的流速都等于绳速，保留动量通量};
\end{tikzpicture}
\figcaption{解图力试2.3A\quad 两直段分别列开系统动量，流出与流入项正确抵消。}
\end{minipage}\end{center}
解图力试2.3A把水平段和斜面段分别列为开系统。斜面段不断接收同速绳元，其沿斜面的动量方程为
\[
\frac{\dd}{\dd t}[\lambda(l-x)v]
=\lambda(l-x)g\sin\theta-T_2+\lambda v^2.
\]
左端展开也有$+\lambda v^2$，抵消后得到
\[
T_1=\lambda x a,\qquad
\lambda(l-x)a=\lambda(l-x)g\sin\theta-T_2.
\]
代入$a$，两端张力相等：
\[
\boxed{T_1=T_2=T
=\frac{\lambda g\sin\theta}{l}x(l-x)}.
\]
在领先阶中，圆弧部分也具有同一张力$T$。对张力曲线求导，$\mathrm dT/\mathrm dx=(\lambda g\sin\theta/l)(l-2x)$，二阶导数为$-2\lambda g\sin\theta/l<0$，确定其峰值位置和下凹形状。下滑的时间方向是$x$从$l$逐渐减小，图C中应沿横轴从右往左读取。
它随$x$先增后减，在$x=l/2$时取最大值$\lambda gl\sin\theta/4$。接触阶段的张力沿这条曲线变化，脱离后的绳形与动力学另由自由运动确定。

\begin{center}\begin{minipage}{.98\linewidth}\centering
\begin{tikzpicture}[scale=.86,transform shape]
\draw[object] plot coordinates {(0.00000,-0.00000) (0.01919,-0.00009) (0.03838,-0.00037) (0.05757,-0.00083) (0.07675,-0.00147) (0.09592,-0.00230) (0.11509,-0.00331) (0.13424,-0.00451) (0.15338,-0.00589) (0.17251,-0.00745) (0.19162,-0.00920) (0.21072,-0.01113) (0.22979,-0.01325) (0.24885,-0.01554) (0.26788,-0.01802) (0.28689,-0.02068) (0.30587,-0.02353) (0.32482,-0.02655) (0.34374,-0.02976) (0.36263,-0.03315) (0.38149,-0.03672) (0.40031,-0.04047) (0.41909,-0.04440) (0.43784,-0.04851) (0.45655,-0.05281) (0.47521,-0.05728) (0.49383,-0.06193) (0.51240,-0.06675) (0.53093,-0.07176) (0.54941,-0.07694) (0.56784,-0.08230) (0.58621,-0.08784) (0.60454,-0.09355) (0.62280,-0.09944) (0.64101,-0.10551) (0.65916,-0.11174) (0.67725,-0.11816) (0.69528,-0.12474) (0.71324,-0.13150) (0.73113,-0.13843) (0.74896,-0.14553) (0.76673,-0.15280) (0.78442,-0.16025) (0.80203,-0.16786) (0.81958,-0.17564) (0.83705,-0.18359) (0.85444,-0.19170) (0.87175,-0.19999) (0.88898,-0.20843) (0.90613,-0.21705) (0.92320,-0.22582) (0.94018,-0.23476) (0.95708,-0.24387) (0.97389,-0.25313) (0.99060,-0.26256) (1.00723,-0.27214) (1.02376,-0.28189) (1.04020,-0.29179) (1.05655,-0.30185) (1.07279,-0.31207) (1.08894,-0.32244) (1.10499,-0.33297) (1.12093,-0.34365) (1.13678,-0.35448) (1.15251,-0.36546) (1.16815,-0.37660) (1.18367,-0.38788) (1.19908,-0.39931) (1.21439,-0.41089) (1.22958,-0.42262) (1.24466,-0.43449) (1.25963,-0.44651) (1.27448,-0.45867) (1.28921,-0.47097) (1.30382,-0.48341) (1.31831,-0.49599) (1.33269,-0.50871) (1.34693,-0.52157) (1.36106,-0.53456) (1.37506,-0.54769) (1.38893,-0.56095) (1.40268,-0.57434) (1.41629,-0.58787) (1.42978,-0.60152) (1.44313,-0.61531) (1.45635,-0.62922) (1.46944,-0.64326) (1.48239,-0.65742) (1.49520,-0.67171) (1.50788,-0.68612) (1.52042,-0.70065) (1.53282,-0.71529) (1.54508,-0.73006) (1.55719,-0.74495) (1.56916,-0.75995) (1.58099,-0.77506) (1.59267,-0.79029) (1.60421,-0.80563) (1.61559,-0.82108) (1.62683,-0.83663)};\draw[object] (-2.00000,0.00000)--(0.00000,0.00000);\draw[guide] (0.00000,-2.00000)--(0.95885,-0.24483);\fill (0.00000,-2.00000) circle(.035);\node[below left] at (0.00000,-2.00000) {曲率中心};\draw[very thick,color=accent] plot coordinates {(0.77884,-0.15788) (0.79152,-0.16329) (0.80417,-0.16880) (0.81678,-0.17438) (0.82935,-0.18006) (0.84188,-0.18582) (0.85437,-0.19167) (0.86682,-0.19761) (0.87923,-0.20363) (0.89160,-0.20974) (0.90393,-0.21593) (0.91621,-0.22220) (0.92845,-0.22856) (0.94064,-0.23501) (0.95279,-0.24154) (0.96490,-0.24815) (0.97696,-0.25485) (0.98897,-0.26163) (1.00093,-0.26849) (1.01285,-0.27543) (1.02472,-0.28246) (1.03654,-0.28957) (1.04831,-0.29676) (1.06003,-0.30403) (1.07171,-0.31138) (1.08333,-0.31881) (1.09489,-0.32632) (1.10641,-0.33391) (1.11787,-0.34158) (1.12928,-0.34933)};\draw[force] (0.77884,-0.15788)--(-0.32644,0.30942);\node[above left] at (-0.32644,0.30942) {$T$};\draw[force] (1.12928,-0.34933)--(2.11969,-1.02690);\node[below right] at (2.11969,-1.02690) {$T$};\draw[force] (0.95885,-0.24483)--(1.67799,1.07154);\node[above right] at (1.67799,1.07154) {$nr\,\dd\phi$};\draw[motion] (0.95885,-0.24483)--(0.33560,-1.38569);\node[below left] at (0.33560,-1.38569) {$v^2/r$向心};\draw[guide] plot coordinates {(0.00000,-1.40000) (0.00882,-1.40006) (0.01764,-1.40026) (0.02646,-1.40058) (0.03527,-1.40104) (0.04408,-1.40162) (0.05287,-1.40233) (0.06166,-1.40318) (0.07043,-1.40415) (0.07918,-1.40525) (0.08792,-1.40648) (0.09664,-1.40783) (0.10533,-1.40932) (0.11401,-1.41093) (0.12266,-1.41267) (0.13128,-1.41454) (0.13988,-1.41653) (0.14844,-1.41865) (0.15698,-1.42090) (0.16547,-1.42327) (0.17394,-1.42577) (0.18236,-1.42838) (0.19075,-1.43113) (0.19909,-1.43399) (0.20740,-1.43698) (0.21565,-1.44009) (0.22386,-1.44333) (0.23202,-1.44668) (0.24014,-1.45015) (0.24820,-1.45374) (0.25620,-1.45745) (0.26415,-1.46128) (0.27205,-1.46522) (0.27988,-1.46928) (0.28766,-1.47345)};\node[left] at (-0.50000,-0.80000) {$\phi$};\draw[guide] (-0.45000,-0.80000)--(0.15000,-1.40000);\node[below] at (1.00000,-3.40000) {领先阶：法向张力合力减向外支持力};
\end{tikzpicture}
\figcaption{解图力试2.3B\quad 弧微元须保留转向加速度，支持方向向外。}
\end{minipage}\end{center}
（2）解图力试2.3B取圆弧上转过角$\phi$、长度$r\dd\phi$的小段。两端张力的合力向曲率中心，大小为$T\dd\phi$；边缘支持力向外，记单位弧长支持力为$n$。该段质量为$\lambda r\dd\phi$，需要的向心加速度为$v^2/r$。于是
\[
T\dd\phi-nr\dd\phi
=\lambda r\dd\phi\,\frac{v^2}{r},
\quad
\boxed{n=\frac{T-\lambda v^2}{r}}.
\]
领先阶的$n$不随$\phi$变化。代入前面的速度和张力，
\[
n=\frac{\lambda g\sin\theta}{rl}(l-x)(2x-l).
\]
支持力密度可以随$r$变小而增大，但积分后的总支持力有限。按图中水平向右、竖直向上的分量，边缘作用方向为$(\sin\phi,\cos\phi)$，故
\[
N_x=\int_0^\theta nr\sin\phi\,\dd\phi
=(T-\lambda v^2)(1-\cos\theta),
\]
\[
N_y=\int_0^\theta nr\cos\phi\,\dd\phi
=(T-\lambda v^2)\sin\theta.
\]
因此边缘对绳的合力大小为
\[
\boxed{N=2(T-\lambda v^2)\sin\frac{\theta}{2}},
\qquad x\ge l/2.
\]
绳对边缘的力与此等大反向。

（3）边缘只能推绳，保持接触时支持密度满足$n\ge0$。除开始零悬长的退化点外，随着$x$减小，首次降到零发生在
\[
\boxed{x_{\rm det}=\frac l2}.
\]
解图力试2.3C比较$T$和$\lambda v^2$：两者在该位置相等。此时张力恰好全部用于提供转向所需的向心力，边缘向外的支持力降到零。贴边公式在$x<l/2$给出负支持力，对应边缘向曲率中心吸引绳。现有凸边只向外支持，故临界点之后绳离开圆弧，后续按自由绳形运动。
\begin{center}\begin{minipage}{.98\linewidth}\centering
\begin{tikzpicture}[scale=.76,transform shape]
\fill[color=accent!7] (3.5,0) rectangle (7,4.3);\draw[axis] (0.00000,0.00000)--(7.45000,0.00000);\draw[axis] (0.00000,0.00000)--(0.00000,4.75000);\node[below right] at (7.45000,0.00000) {$x/l$};\node[above right] at (0.00000,4.75000) {$4F/(\lambda gl\sin\theta)$};\draw[object] (-0.07000,1.00000)--(0.07000,1.00000);\node[left] at (0.00000,1.00000) {1};\draw[object] (-0.07000,4.00000)--(0.07000,4.00000);\node[left] at (0.00000,4.00000) {4};\draw[traj,dashed] plot coordinates {(0.00000,0.00000) (0.03535,0.02010) (0.07071,0.04000) (0.10606,0.05969) (0.14141,0.07918) (0.17677,0.09846) (0.21212,0.11754) (0.24747,0.13641) (0.28283,0.15509) (0.31818,0.17355) (0.35354,0.19182) (0.38889,0.20988) (0.42424,0.22773) (0.45960,0.24538) (0.49495,0.26283) (0.53030,0.28007) (0.56566,0.29711) (0.60101,0.31395) (0.63636,0.33058) (0.67172,0.34701) (0.70707,0.36323) (0.74242,0.37925) (0.77778,0.39506) (0.81313,0.41067) (0.84848,0.42608) (0.88384,0.44128) (0.91919,0.45628) (0.95455,0.47107) (0.98990,0.48566) (1.02525,0.50005) (1.06061,0.51423) (1.09596,0.52821) (1.13131,0.54199) (1.16667,0.55556) (1.20202,0.56892) (1.23737,0.58208) (1.27273,0.59504) (1.30808,0.60780) (1.34343,0.62034) (1.37879,0.63269) (1.41414,0.64483) (1.44949,0.65677) (1.48485,0.66850) (1.52020,0.68003) (1.55556,0.69136) (1.59091,0.70248) (1.62626,0.71340) (1.66162,0.72411) (1.69697,0.73462) (1.73232,0.74492) (1.76768,0.75502) (1.80303,0.76492) (1.83838,0.77461) (1.87374,0.78410) (1.90909,0.79339) (1.94444,0.80247) (1.97980,0.81135) (2.01515,0.82002) (2.05051,0.82849) (2.08586,0.83675) (2.12121,0.84481) (2.15657,0.85267) (2.19192,0.86032) (2.22727,0.86777) (2.26263,0.87501) (2.29798,0.88205) (2.33333,0.88889) (2.36869,0.89552) (2.40404,0.90195) (2.43939,0.90817) (2.47475,0.91419) (2.51010,0.92001) (2.54545,0.92562) (2.58081,0.93103) (2.61616,0.93623) (2.65152,0.94123) (2.68687,0.94603) (2.72222,0.95062) (2.75758,0.95500) (2.79293,0.95919) (2.82828,0.96317) (2.86364,0.96694) (2.89899,0.97051) (2.93434,0.97388) (2.96970,0.97704) (3.00505,0.98000) (3.04040,0.98276) (3.07576,0.98531) (3.11111,0.98765) (3.14646,0.98980) (3.18182,0.99174) (3.21717,0.99347) (3.25253,0.99500) (3.28788,0.99633) (3.32323,0.99745) (3.35859,0.99837) (3.39394,0.99908) (3.42929,0.99959) (3.46465,0.99990) (3.50000,1.00000)};\draw[force,dashed] plot coordinates {(0.00000,4.00000) (0.03535,3.95970) (0.07071,3.91960) (0.10606,3.87971) (0.14141,3.84002) (0.17677,3.80053) (0.21212,3.76125) (0.24747,3.72217) (0.28283,3.68330) (0.31818,3.64463) (0.35354,3.60616) (0.38889,3.56790) (0.42424,3.52984) (0.45960,3.49199) (0.49495,3.45434) (0.53030,3.41690) (0.56566,3.37966) (0.60101,3.34262) (0.63636,3.30579) (0.67172,3.26916) (0.70707,3.23273) (0.74242,3.19651) (0.77778,3.16049) (0.81313,3.12468) (0.84848,3.08907) (0.88384,3.05367) (0.91919,3.01847) (0.95455,2.98347) (0.98990,2.94868) (1.02525,2.91409) (1.06061,2.87971) (1.09596,2.84553) (1.13131,2.81155) (1.16667,2.77778) (1.20202,2.74421) (1.23737,2.71085) (1.27273,2.67769) (1.30808,2.64473) (1.34343,2.61198) (1.37879,2.57943) (1.41414,2.54709) (1.44949,2.51495) (1.48485,2.48301) (1.52020,2.45128) (1.55556,2.41975) (1.59091,2.38843) (1.62626,2.35731) (1.66162,2.32640) (1.69697,2.29568) (1.73232,2.26518) (1.76768,2.23487) (1.80303,2.20478) (1.83838,2.17488) (1.87374,2.14519) (1.90909,2.11570) (1.94444,2.08642) (1.97980,2.05734) (2.01515,2.02847) (2.05051,1.99980) (2.08586,1.97133) (2.12121,1.94307) (2.15657,1.91501) (2.19192,1.88715) (2.22727,1.85950) (2.26263,1.83206) (2.29798,1.80482) (2.33333,1.77778) (2.36869,1.75094) (2.40404,1.72431) (2.43939,1.69789) (2.47475,1.67167) (2.51010,1.64565) (2.54545,1.61983) (2.58081,1.59423) (2.61616,1.56882) (2.65152,1.54362) (2.68687,1.51862) (2.72222,1.49383) (2.75758,1.46924) (2.79293,1.44485) (2.82828,1.42067) (2.86364,1.39669) (2.89899,1.37292) (2.93434,1.34935) (2.96970,1.32599) (3.00505,1.30283) (3.04040,1.27987) (3.07576,1.25712) (3.11111,1.23457) (3.14646,1.21222) (3.18182,1.19008) (3.21717,1.16815) (3.25253,1.14641) (3.28788,1.12489) (3.32323,1.10356) (3.35859,1.08244) (3.39394,1.06152) (3.42929,1.04081) (3.46465,1.02030) (3.50000,1.00000)};\draw[traj] plot coordinates {(3.50000,1.00000) (3.53535,0.99990) (3.57071,0.99959) (3.60606,0.99908) (3.64141,0.99837) (3.67677,0.99745) (3.71212,0.99633) (3.74747,0.99500) (3.78283,0.99347) (3.81818,0.99174) (3.85354,0.98980) (3.88889,0.98765) (3.92424,0.98531) (3.95960,0.98276) (3.99495,0.98000) (4.03030,0.97704) (4.06566,0.97388) (4.10101,0.97051) (4.13636,0.96694) (4.17172,0.96317) (4.20707,0.95919) (4.24242,0.95500) (4.27778,0.95062) (4.31313,0.94603) (4.34848,0.94123) (4.38384,0.93623) (4.41919,0.93103) (4.45455,0.92562) (4.48990,0.92001) (4.52525,0.91419) (4.56061,0.90817) (4.59596,0.90195) (4.63131,0.89552) (4.66667,0.88889) (4.70202,0.88205) (4.73737,0.87501) (4.77273,0.86777) (4.80808,0.86032) (4.84343,0.85267) (4.87879,0.84481) (4.91414,0.83675) (4.94949,0.82849) (4.98485,0.82002) (5.02020,0.81135) (5.05556,0.80247) (5.09091,0.79339) (5.12626,0.78410) (5.16162,0.77461) (5.19697,0.76492) (5.23232,0.75502) (5.26768,0.74492) (5.30303,0.73462) (5.33838,0.72411) (5.37374,0.71340) (5.40909,0.70248) (5.44444,0.69136) (5.47980,0.68003) (5.51515,0.66850) (5.55051,0.65677) (5.58586,0.64483) (5.62121,0.63269) (5.65657,0.62034) (5.69192,0.60780) (5.72727,0.59504) (5.76263,0.58208) (5.79798,0.56892) (5.83333,0.55556) (5.86869,0.54199) (5.90404,0.52821) (5.93939,0.51423) (5.97475,0.50005) (6.01010,0.48566) (6.04545,0.47107) (6.08081,0.45628) (6.11616,0.44128) (6.15152,0.42608) (6.18687,0.41067) (6.22222,0.39506) (6.25758,0.37925) (6.29293,0.36323) (6.32828,0.34701) (6.36364,0.33058) (6.39899,0.31395) (6.43434,0.29711) (6.46970,0.28007) (6.50505,0.26283) (6.54040,0.24538) (6.57576,0.22773) (6.61111,0.20988) (6.64646,0.19182) (6.68182,0.17355) (6.71717,0.15509) (6.75253,0.13641) (6.78788,0.11754) (6.82323,0.09846) (6.85859,0.07918) (6.89394,0.05969) (6.92929,0.04000) (6.96465,0.02010) (7.00000,0.00000)};\draw[force] plot coordinates {(3.50000,1.00000) (3.53535,0.97990) (3.57071,0.96000) (3.60606,0.94031) (3.64141,0.92082) (3.67677,0.90154) (3.71212,0.88246) (3.74747,0.86359) (3.78283,0.84491) (3.81818,0.82645) (3.85354,0.80818) (3.88889,0.79012) (3.92424,0.77227) (3.95960,0.75462) (3.99495,0.73717) (4.03030,0.71993) (4.06566,0.70289) (4.10101,0.68605) (4.13636,0.66942) (4.17172,0.65299) (4.20707,0.63677) (4.24242,0.62075) (4.27778,0.60494) (4.31313,0.58933) (4.34848,0.57392) (4.38384,0.55872) (4.41919,0.54372) (4.45455,0.52893) (4.48990,0.51434) (4.52525,0.49995) (4.56061,0.48577) (4.59596,0.47179) (4.63131,0.45801) (4.66667,0.44444) (4.70202,0.43108) (4.73737,0.41792) (4.77273,0.40496) (4.80808,0.39220) (4.84343,0.37966) (4.87879,0.36731) (4.91414,0.35517) (4.94949,0.34323) (4.98485,0.33150) (5.02020,0.31997) (5.05556,0.30864) (5.09091,0.29752) (5.12626,0.28660) (5.16162,0.27589) (5.19697,0.26538) (5.23232,0.25508) (5.26768,0.24498) (5.30303,0.23508) (5.33838,0.22539) (5.37374,0.21590) (5.40909,0.20661) (5.44444,0.19753) (5.47980,0.18865) (5.51515,0.17998) (5.55051,0.17151) (5.58586,0.16325) (5.62121,0.15519) (5.65657,0.14733) (5.69192,0.13968) (5.72727,0.13223) (5.76263,0.12499) (5.79798,0.11795) (5.83333,0.11111) (5.86869,0.10448) (5.90404,0.09805) (5.93939,0.09183) (5.97475,0.08581) (6.01010,0.07999) (6.04545,0.07438) (6.08081,0.06897) (6.11616,0.06377) (6.15152,0.05877) (6.18687,0.05397) (6.22222,0.04938) (6.25758,0.04500) (6.29293,0.04081) (6.32828,0.03683) (6.36364,0.03306) (6.39899,0.02949) (6.43434,0.02612) (6.46970,0.02296) (6.50505,0.02000) (6.54040,0.01724) (6.57576,0.01469) (6.61111,0.01235) (6.64646,0.01020) (6.68182,0.00826) (6.71717,0.00653) (6.75253,0.00500) (6.78788,0.00367) (6.82323,0.00255) (6.85859,0.00163) (6.89394,0.00092) (6.92929,0.00041) (6.96465,0.00010) (7.00000,0.00000)};\draw[guide] (3.50000,0.00000)--(3.50000,4.30000);\fill (3.50000,1.00000) circle(.035);\draw[object] (3.50000,-0.07000)--(3.50000,0.07000);\node[below] at (3.50000,-0.10000) {1/2};\draw[object] (7.00000,-0.07000)--(7.00000,0.07000);\node[below] at (7.00000,-0.10000) {1};\draw[traj] (7.80000,4.10000)--(8.50000,4.10000);\node[right] at (8.70000,4.10000) {$T$};\draw[force] (7.80000,3.50000)--(8.50000,3.50000);\node[right] at (8.70000,3.50000) {$\lambda v^2$};\node[anchor=center] at (5.25000,3.60000) {保持接触};\draw[motion] (7.00000,-1.00000)--(3.50000,-1.00000);\node[below] at (5.25000,-1.10000) {下滑从右向左};\node[below] at (3.80000,-2.50000) {虚线为贴边方程的形式延伸；左侧需要负支持力};
\end{tikzpicture}
\figcaption{解图力试2.3C\quad 接触要求$T\ge\lambda v^2$，在$x/l=1/2$开始失效。}
\end{minipage}\end{center}

（4）有限$r$的修正同时计入圆弧长度、重力和沿弧张力变化。圆弧被绳完整占据时，设斜面直段长度
\[
y=l-x-r\theta.
\]
圆弧的顶切点高度取0，弧上高度为$-r(1-\cos\phi)$，底切点高度为$-r(1-\cos\theta)$。计入圆弧和斜段势能，能量与加速度分别为
\[
v^2=\frac gl\left[
2r^2(\theta-\sin\theta)+2ry(1-\cos\theta)+y^2\sin\theta
\right],
\]
\[
a=\frac gl\left[r(1-\cos\theta)+y\sin\theta\right].
\]
圆弧上切向牛顿方程为$\dd T/\dd\phi=\lambda r(a-g\sin\phi)$，顶端$T(0)=\lambda xa$，故
\[
T(\phi)=\lambda xa+\lambda r\left[a\phi+g(\cos\phi-1)\right].
\]
法向方程还要计入重力的向内分量：
\[
\boxed{n(\phi)=\frac{T(\phi)-\lambda v^2}{r}
+\lambda g\cos\phi}.
\]
严格的首次脱离条件是$\min_{0\le\phi\le\theta}n(\phi)=0$，张力也不再沿弧完全相同。让$r/l\to0$，这些式子恢复$x_{\rm det}/l\to1/2$；有限半径的脱离位置由实际角分布的最小支持密度确定，$x=l/2$是$r/l\to0$时的极限。

\begin{center}\begin{minipage}{.98\linewidth}\centering
\begin{tikzpicture}[scale=.8,transform shape]
\draw[axis] (0,0)--(10.7,0) node[right] {$\phi$};
\draw[axis] (0,0)--(0,3.2) node[above] {$n/(\lambda g)$};
\node[above] at (5,3.5) {$\theta=60^\circ,r/l=0.02$：$x/l=0.55,0.50$及首次脱离的角分布};\draw[guide] (.2,-1)--(.9,-1);\node[right] at (1,-1) {$x/l=0.55$};\draw[traj] (.2,-1.5)--(.9,-1.5);\node[right] at (1,-1.5) {$x/l=0.50$};\draw[dashed,thick,color=warm] (.2,-2)--(.9,-2);\node[right] at (1,-2) {首次脱离，$x/l\simeq0.484863$};
\draw[guide] plot coordinates {(0.000000000,1.542837525) (0.083333333,1.544464376) (0.166666667,1.546015077) (0.250000000,1.547489635) (0.333333333,1.548888065) (0.416666667,1.550210388) (0.500000000,1.551456629) (0.583333333,1.552626821) (0.666666667,1.553721001) (0.750000000,1.554739213) (0.833333333,1.555681505) (0.916666667,1.556547934) (1.000000000,1.557338559) (1.083333333,1.558053447) (1.166666667,1.558692671) (1.250000000,1.559256309) (1.333333333,1.559744445) (1.416666667,1.560157168) (1.500000000,1.560494573) (1.583333333,1.560756762) (1.666666667,1.560943842) (1.750000000,1.561055925) (1.833333333,1.561093129) (1.916666667,1.561055578) (2.000000000,1.560943403) (2.083333333,1.560756737) (2.166666667,1.560495723) (2.250000000,1.560160507) (2.333333333,1.559751241) (2.416666667,1.559268083) (2.500000000,1.558711197) (2.583333333,1.558080753) (2.666666667,1.557376924) (2.750000000,1.556599891) (2.833333333,1.555749840) (2.916666667,1.554826963) (3.000000000,1.553831457) (3.083333333,1.552763524) (3.166666667,1.551623372) (3.250000000,1.550411216) (3.333333333,1.549127274) (3.416666667,1.547771771) (3.500000000,1.546344936) (3.583333333,1.544847006) (3.666666667,1.543278221) (3.750000000,1.541638827) (3.833333333,1.539929076) (3.916666667,1.538149225) (4.000000000,1.536299537) (4.083333333,1.534380278) (4.166666667,1.532391722) (4.250000000,1.530334148) (4.333333333,1.528207838) (4.416666667,1.526013082) (4.500000000,1.523750172) (4.583333333,1.521419410) (4.666666667,1.519021097) (4.750000000,1.516555545) (4.833333333,1.514023068) (4.916666667,1.511423985) (5.000000000,1.508758621) (5.083333333,1.506027306) (5.166666667,1.503230375) (5.250000000,1.500368166) (5.333333333,1.497441026) (5.416666667,1.494449304) (5.500000000,1.491393355) (5.583333333,1.488273537) (5.666666667,1.485090216) (5.750000000,1.481843760) (5.833333333,1.478534544) (5.916666667,1.475162946) (6.000000000,1.471729350) (6.083333333,1.468234145) (6.166666667,1.464677722) (6.250000000,1.461060481) (6.333333333,1.457382822) (6.416666667,1.453645154) (6.500000000,1.449847887) (6.583333333,1.445991437) (6.666666667,1.442076225) (6.750000000,1.438102675) (6.833333333,1.434071218) (6.916666667,1.429982287) (7.000000000,1.425836320) (7.083333333,1.421633759) (7.166666667,1.417375052) (7.250000000,1.413060650) (7.333333333,1.408691008) (7.416666667,1.404266585) (7.500000000,1.399787845) (7.583333333,1.395255256) (7.666666667,1.390669291) (7.750000000,1.386030424) (7.833333333,1.381339137) (7.916666667,1.376595912) (8.000000000,1.371801239) (8.083333333,1.366955609) (8.166666667,1.362059518) (8.250000000,1.357113466) (8.333333333,1.352117955) (8.416666667,1.347073494) (8.500000000,1.341980593) (8.583333333,1.336839767) (8.666666667,1.331651534) (8.750000000,1.326416416) (8.833333333,1.321134938) (8.916666667,1.315807630) (9.000000000,1.310435024) (9.083333333,1.305017655) (9.166666667,1.299556064) (9.250000000,1.294050793) (9.333333333,1.288502388) (9.416666667,1.282911398) (9.500000000,1.277278376) (9.583333333,1.271603877) (9.666666667,1.265888461) (9.750000000,1.260132690) (9.833333333,1.254337128) (9.916666667,1.248502345) (10.000000000,1.242628910)};
\draw[traj] plot coordinates {(0.000000000,0.599076422) (0.083333333,0.600892211) (0.166666667,0.602631848) (0.250000000,0.604295344) (0.333333333,0.605882711) (0.416666667,0.607393972) (0.500000000,0.608829150) (0.583333333,0.610188280) (0.666666667,0.611471397) (0.750000000,0.612678546) (0.833333333,0.613809776) (0.916666667,0.614865142) (1.000000000,0.615844705) (1.083333333,0.616748531) (1.166666667,0.617576692) (1.250000000,0.618329268) (1.333333333,0.619006341) (1.416666667,0.619608001) (1.500000000,0.620134344) (1.583333333,0.620585471) (1.666666667,0.620961488) (1.750000000,0.621262508) (1.833333333,0.621488649) (1.916666667,0.621640036) (2.000000000,0.621716798) (2.083333333,0.621719070) (2.166666667,0.621646993) (2.250000000,0.621500715) (2.333333333,0.621280386) (2.416666667,0.620986166) (2.500000000,0.620618217) (2.583333333,0.620176710) (2.666666667,0.619661818) (2.750000000,0.619073723) (2.833333333,0.618412610) (2.916666667,0.617678670) (3.000000000,0.616872101) (3.083333333,0.615993106) (3.166666667,0.615041892) (3.250000000,0.614018673) (3.333333333,0.612923668) (3.416666667,0.611757102) (3.500000000,0.610519205) (3.583333333,0.609210212) (3.666666667,0.607830365) (3.750000000,0.606379908) (3.833333333,0.604859095) (3.916666667,0.603268181) (4.000000000,0.601607430) (4.083333333,0.599877109) (4.166666667,0.598077491) (4.250000000,0.596208854) (4.333333333,0.594271482) (4.416666667,0.592265662) (4.500000000,0.590191691) (4.583333333,0.588049865) (4.666666667,0.585840491) (4.750000000,0.583563876) (4.833333333,0.581220336) (4.916666667,0.578810191) (5.000000000,0.576333764) (5.083333333,0.573791386) (5.166666667,0.571183392) (5.250000000,0.568510122) (5.333333333,0.565771919) (5.416666667,0.562969134) (5.500000000,0.560102122) (5.583333333,0.557171242) (5.666666667,0.554176858) (5.750000000,0.551119340) (5.833333333,0.547999061) (5.916666667,0.544816401) (6.000000000,0.541571742) (6.083333333,0.538265474) (6.166666667,0.534897989) (6.250000000,0.531469685) (6.333333333,0.527980964) (6.416666667,0.524432233) (6.500000000,0.520823903) (6.583333333,0.517156391) (6.666666667,0.513430116) (6.750000000,0.509645504) (6.833333333,0.505802985) (6.916666667,0.501902991) (7.000000000,0.497945961) (7.083333333,0.493932338) (7.166666667,0.489862569) (7.250000000,0.485737104) (7.333333333,0.481556399) (7.416666667,0.477320913) (7.500000000,0.473031111) (7.583333333,0.468687460) (7.666666667,0.464290431) (7.750000000,0.459840502) (7.833333333,0.455338152) (7.916666667,0.450783865) (8.000000000,0.446178130) (8.083333333,0.441521437) (8.166666667,0.436814284) (8.250000000,0.432057169) (8.333333333,0.427250596) (8.416666667,0.422395072) (8.500000000,0.417491108) (8.583333333,0.412539220) (8.666666667,0.407539924) (8.750000000,0.402493743) (8.833333333,0.397401203) (8.916666667,0.392262832) (9.000000000,0.387079163) (9.083333333,0.381850733) (9.166666667,0.376578079) (9.250000000,0.371261745) (9.333333333,0.365902277) (9.416666667,0.360500225) (9.500000000,0.355056140) (9.583333333,0.349570579) (9.666666667,0.344044101) (9.750000000,0.338477267) (9.833333333,0.332870642) (9.916666667,0.327224796) (10.000000000,0.321540299)};
\draw[dashed,thick,color=warm] plot coordinates {(0.000000000,0.270672340) (0.083333333,0.272545327) (0.166666667,0.274342163) (0.250000000,0.276062856) (0.333333333,0.277707422) (0.416666667,0.279275880) (0.500000000,0.280768257) (0.583333333,0.282184585) (0.666666667,0.283524901) (0.750000000,0.284789248) (0.833333333,0.285977676) (0.916666667,0.287090240) (1.000000000,0.288127001) (1.083333333,0.289088025) (1.166666667,0.289973385) (1.250000000,0.290783159) (1.333333333,0.291517430) (1.416666667,0.292176288) (1.500000000,0.292759829) (1.583333333,0.293268154) (1.666666667,0.293701369) (1.750000000,0.294059588) (1.833333333,0.294342928) (1.916666667,0.294551513) (2.000000000,0.294685473) (2.083333333,0.294744943) (2.166666667,0.294730064) (2.250000000,0.294640984) (2.333333333,0.294477853) (2.416666667,0.294240831) (2.500000000,0.293930081) (2.583333333,0.293545772) (2.666666667,0.293088078) (2.750000000,0.292557181) (2.833333333,0.291953266) (2.916666667,0.291276525) (3.000000000,0.290527154) (3.083333333,0.289705357) (3.166666667,0.288811341) (3.250000000,0.287845320) (3.333333333,0.286807514) (3.416666667,0.285698146) (3.500000000,0.284517447) (3.583333333,0.283265653) (3.666666667,0.281943003) (3.750000000,0.280549745) (3.833333333,0.279086130) (3.916666667,0.277552414) (4.000000000,0.275948861) (4.083333333,0.274275739) (4.166666667,0.272533319) (4.250000000,0.270721880) (4.333333333,0.268841705) (4.416666667,0.266893085) (4.500000000,0.264876311) (4.583333333,0.262791684) (4.666666667,0.260639507) (4.750000000,0.258420091) (4.833333333,0.256133749) (4.916666667,0.253780802) (5.000000000,0.251361574) (5.083333333,0.248876394) (5.166666667,0.246325598) (5.250000000,0.243709526) (5.333333333,0.241028521) (5.416666667,0.238282935) (5.500000000,0.235473121) (5.583333333,0.232599439) (5.666666667,0.229662253) (5.750000000,0.226661933) (5.833333333,0.223598853) (5.916666667,0.220473390) (6.000000000,0.217285930) (6.083333333,0.214036860) (6.166666667,0.210726574) (6.250000000,0.207355468) (6.333333333,0.203923945) (6.416666667,0.200432412) (6.500000000,0.196881280) (6.583333333,0.193270966) (6.666666667,0.189601890) (6.750000000,0.185874476) (6.833333333,0.182089154) (6.916666667,0.178246359) (7.000000000,0.174346527) (7.083333333,0.170390103) (7.166666667,0.166377531) (7.250000000,0.162309264) (7.333333333,0.158185758) (7.416666667,0.154007470) (7.500000000,0.149774866) (7.583333333,0.145488413) (7.666666667,0.141148583) (7.750000000,0.136755852) (7.833333333,0.132310700) (7.916666667,0.127813612) (8.000000000,0.123265074) (8.083333333,0.118665580) (8.166666667,0.114015625) (8.250000000,0.109315708) (8.333333333,0.104566333) (8.416666667,0.099768007) (8.500000000,0.094921242) (8.583333333,0.090026551) (8.666666667,0.085084454) (8.750000000,0.080095471) (8.833333333,0.075060129) (8.916666667,0.069978957) (9.000000000,0.064852486) (9.083333333,0.059681254) (9.166666667,0.054465798) (9.250000000,0.049206662) (9.333333333,0.043904393) (9.416666667,0.038559538) (9.500000000,0.033172652) (9.583333333,0.027744289) (9.666666667,0.022275009) (9.750000000,0.016765373) (9.833333333,0.011215947) (9.916666667,0.005627299) (10.000000000,0.000000000)};
\node[below] at (0,-.12) {$0$};
\node[below] at (5,-.12) {$30^\circ$};
\node[below] at (10,-.12) {$60^\circ$};
\end{tikzpicture}
\figcaption{解图力试2.3D\quad $\theta=60^\circ,r/l=0.02$：$x/l=0.55,0.50$及首次脱离的角分布}
\end{minipage}\end{center}
图D按完整圆弧的法向式计算，靠近底切点的支持最小，首次脱离时该端点降至零，其余弧段仍保持正支持。有限半径的临界张力同时与速度项、重力法向分量平衡。
初始完全水平、零悬长且严格静止的理想绳对应一个静止解。本题的微扰用于选出下滑过程；上式从微小下垂段的能量趋零极限建立，脱离位置的极限与由微扰决定的启动时间分别计算。


\Needspace{13\baselineskip}

\pdfbookmark[2]{第 4 题 两细杆端部弹性碰撞与接触点速度}{q4}

\section*{第 4 题\quad 两细杆端部弹性碰撞与接触点速度}\label{q:4}


光滑水平桌面上，两根均匀细杆$A,B$平行于$x$轴，长度均为$l$，质量分别为$m_A,m_B$。杆$A$静止，位于$x$轴的$(-l+\epsilon,\epsilon)$区域，$0<\epsilon\ll l$。杆$B$位于其下方，横坐标覆盖$(0,l)$，以$v_0\hat y$作无转动平移。

当杆$B$向上到达$x$轴时，其左端与杆$A$右端发生光滑、单点、完全弹性碰撞。取$\epsilon/l\to0$的极限，接触法向沿$y$；忽略杆厚和接触时间内的位移。
\begin{enumerate}
\item 求碰后两杆质心的带符号速度$v_A,v_B$及角速度$\omega_A,\omega_B$，其中向$+y$、逆时针分别为正。
\item 验证
\[
v_0=\left(v_A+\frac l2\omega_A\right)
-\left(v_B-\frac l2\omega_B\right),
\]
并说明它在本题中的物理依据和适用条件。
\end{enumerate}


\begin{center}\begin{minipage}{.98\linewidth}\centering
\begin{tikzpicture}[scale=.9,transform shape]
\draw[axis] (0.00000,0.00000)--(5.25000,0.00000);\draw[axis] (0.00000,-2.20000)--(0.00000,2.25000);\node[below right] at (5.25000,0.00000) {$x$};\node[above right] at (0.00000,2.25000) {$y$};\draw[very thick,color=accent] (-4.00000,0.00000)--(0.00000,0.00000);\draw[very thick,color=accent] (0.00000,-1.50000)--(4.00000,-1.50000);\node[below] at (-2.00000,-0.10000) {杆 $A,m_A$};\node[below] at (2.00000,-1.60000) {杆 $B,m_B$};\draw[motion] (2.00000,-1.40000)--(2.00000,-0.30000);\node[right] at (2.10000,-0.70000) {$v_0$};\fill (0.00000,0.00000) circle(.035);\node[below left] at (0.00000,0.00000) {$O$};
\end{tikzpicture}
\figcaption{图力试2.4\quad B杆向+y平移，端部将与静止A杆接触，取小重叠极限。}
\end{minipage}\end{center}

\Needspace{6\baselineskip}\subsection*{解答}


图力试2.4中的接触冲量作用在两杆相对的端部。设杆$A$受到$+J\hat y$，杆$B$受到$-J\hat y$，其中$J>0$。

\begin{center}\begin{minipage}{.98\linewidth}\centering
\begin{tikzpicture}[scale=.82,transform shape]
\draw[very thick,color=accent] (-4.00000,0.00000)--(0.00000,0.00000);\draw[very thick,color=ink] (0.00000,0.00000)--(4.00000,0.00000);\fill (-2.00000,0.00000) circle(.035);\node[below] at (-2.00000,0.00000) {$C_A$};\fill (2.00000,0.00000) circle(.035);\node[below] at (2.00000,0.00000) {$C_B$};\draw[force] (-0.12000,0.00000)--(-0.12000,1.70000);\node[above left] at (-0.20000,1.70000) {$+J$};\draw[force] (0.12000,0.00000)--(0.12000,-1.70000);\node[below right] at (0.20000,-1.70000) {$-J$};\draw[motion] plot coordinates {(-1.15425,0.08486) (-1.15840,0.11924) (-1.16396,0.15341) (-1.17090,0.18734) (-1.17922,0.22095) (-1.18890,0.25420) (-1.19993,0.28702) (-1.21228,0.31937) (-1.22594,0.35119) (-1.24089,0.38243) (-1.25709,0.41303) (-1.27453,0.44294) (-1.29318,0.47212) (-1.31299,0.50052) (-1.33395,0.52808) (-1.35601,0.55477) (-1.37914,0.58054) (-1.40330,0.60535) (-1.42845,0.62915) (-1.45455,0.65191) (-1.48155,0.67358) (-1.50942,0.69414) (-1.53810,0.71354) (-1.56754,0.73177) (-1.59771,0.74877) (-1.62854,0.76454) (-1.65998,0.77903) (-1.69200,0.79223) (-1.72452,0.80412) (-1.75750,0.81467) (-1.79088,0.82387) (-1.82461,0.83171) (-1.85863,0.83816) (-1.89289,0.84322) (-1.92732,0.84689) (-1.96187,0.84914) (-1.99649,0.84999) (-2.03111,0.84943) (-2.06568,0.84746) (-2.10015,0.84408) (-2.13444,0.83930) (-2.16851,0.83313) (-2.20231,0.82557) (-2.23576,0.81665) (-2.26883,0.80637) (-2.30145,0.79475) (-2.33357,0.78181) (-2.36514,0.76758) (-2.39610,0.75207) (-2.42640,0.73531) (-2.45599,0.71734) (-2.48483,0.69817) (-2.51287,0.67784) (-2.54005,0.65639) (-2.56633,0.63385)};\node[above] at (-2.00000,1.30000) {$\omega_A>0$};\draw[motion] plot coordinates {(2.84575,0.08486) (2.84160,0.11924) (2.83604,0.15341) (2.82910,0.18734) (2.82078,0.22095) (2.81110,0.25420) (2.80007,0.28702) (2.78772,0.31937) (2.77406,0.35119) (2.75911,0.38243) (2.74291,0.41303) (2.72547,0.44294) (2.70682,0.47212) (2.68701,0.50052) (2.66605,0.52808) (2.64399,0.55477) (2.62086,0.58054) (2.59670,0.60535) (2.57155,0.62915) (2.54545,0.65191) (2.51845,0.67358) (2.49058,0.69414) (2.46190,0.71354) (2.43246,0.73177) (2.40229,0.74877) (2.37146,0.76454) (2.34002,0.77903) (2.30800,0.79223) (2.27548,0.80412) (2.24250,0.81467) (2.20912,0.82387) (2.17539,0.83171) (2.14137,0.83816) (2.10711,0.84322) (2.07268,0.84689) (2.03813,0.84914) (2.00351,0.84999) (1.96889,0.84943) (1.93432,0.84746) (1.89985,0.84408) (1.86556,0.83930) (1.83149,0.83313) (1.79769,0.82557) (1.76424,0.81665) (1.73117,0.80637) (1.69855,0.79475) (1.66643,0.78181) (1.63486,0.76758) (1.60390,0.75207) (1.57360,0.73531) (1.54401,0.71734) (1.51517,0.69817) (1.48713,0.67784) (1.45995,0.65639) (1.43367,0.63385)};\node[above] at (2.00000,1.30000) {$\omega_B>0$};\draw[<->,thin,color=muted] (-2.00000,-2.50000)--(0.00000,-2.50000);\node[below] at (-1.00000,-2.50000) {$l/2$};\draw[<->,thin,color=muted] (0.00000,-3.50000)--(2.00000,-3.50000);\node[below] at (1.00000,-3.50000) {$l/2$};\node[above] at (0.00000,2.80000) {冲量相反，绕各自质心的力矩同向};
\end{tikzpicture}
\figcaption{解图力试2.4A\quad 两端接触冲量关于杆质心产生同向角冲量。}
\end{minipage}\end{center}
解图力试2.4A标明关于各自质心的力臂：杆$A$在右端受向上冲量，杆$B$在左端受向下冲量，两者都产生逆时针角冲量。由平动与转动冲量方程，
\[
v_A=\frac{J}{m_A},\qquad
v_B=v_0-\frac{J}{m_B},
\]
\[
I_A\omega_A=\frac{Jl}{2},\quad
I_B\omega_B=\frac{Jl}{2},\quad
I_A=\frac{m_Al^2}{12},\ I_B=\frac{m_Bl^2}{12}.
\]
因此
\[
\omega_A=\frac{6J}{m_Al},\qquad
\omega_B=\frac{6J}{m_Bl}.
\]
用同一个接触冲量表示四个速度后，总动能就成为冲量的一元二次函数；逐杆角冲量同时给出两杆角速度的联系。

（1）完全弹性刚体模型要求碰前、碰后的总平动加转动动能相同。把上述速度代入，动能改变量为
\[
\Delta K=-Jv_0+
\frac{J^2}{2}
\left(\frac1{m_A}+\frac1{m_B}
+\frac{(l/2)^2}{I_A}+\frac{(l/2)^2}{I_B}\right).
\]
其中两项$1/m$来自质心平动，两项$(l/2)^2/I$来自绕质心转动。每项转动响应等于$3/m$，所以转动能与平动能共同决定接触的惯性。记接触点的有效逆质量
\[
\mathcal K=
\frac1{m_A}+\frac1{m_B}
+\frac{(l/2)^2}{I_A}+\frac{(l/2)^2}{I_B}
=4\left(\frac1{m_A}+\frac1{m_B}\right).
\]
\begin{center}\begin{minipage}{.98\linewidth}\centering
\begin{tikzpicture}[scale=.8,transform shape]
\draw[axis] (0,0)--(8.5,0) node[right] {$J/J_{\rm el}$};
\draw[axis] (0,0)--(0,3.2) node[above] {$\mathcal K\Delta K/v_0^2$};
\node[above] at (5,3.5) {累计冲量增大：总动能先降后回到初值};
\draw[traj] plot coordinates {(0.000000000,2.000000000) (0.080000000,1.920800000) (0.160000000,1.843200000) (0.240000000,1.767200000) (0.320000000,1.692800000) (0.400000000,1.620000000) (0.480000000,1.548800000) (0.560000000,1.479200000) (0.640000000,1.411200000) (0.720000000,1.344800000) (0.800000000,1.280000000) (0.880000000,1.216800000) (0.960000000,1.155200000) (1.040000000,1.095200000) (1.120000000,1.036800000) (1.200000000,0.980000000) (1.280000000,0.924800000) (1.360000000,0.871200000) (1.440000000,0.819200000) (1.520000000,0.768800000) (1.600000000,0.720000000) (1.680000000,0.672800000) (1.760000000,0.627200000) (1.840000000,0.583200000) (1.920000000,0.540800000) (2.000000000,0.500000000) (2.080000000,0.460800000) (2.160000000,0.423200000) (2.240000000,0.387200000) (2.320000000,0.352800000) (2.400000000,0.320000000) (2.480000000,0.288800000) (2.560000000,0.259200000) (2.640000000,0.231200000) (2.720000000,0.204800000) (2.800000000,0.180000000) (2.880000000,0.156800000) (2.960000000,0.135200000) (3.040000000,0.115200000) (3.120000000,0.096800000) (3.200000000,0.080000000) (3.280000000,0.064800000) (3.360000000,0.051200000) (3.440000000,0.039200000) (3.520000000,0.028800000) (3.600000000,0.020000000) (3.680000000,0.012800000) (3.760000000,0.007200000) (3.840000000,0.003200000) (3.920000000,0.000800000) (4.000000000,0.000000000) (4.080000000,0.000800000) (4.160000000,0.003200000) (4.240000000,0.007200000) (4.320000000,0.012800000) (4.400000000,0.020000000) (4.480000000,0.028800000) (4.560000000,0.039200000) (4.640000000,0.051200000) (4.720000000,0.064800000) (4.800000000,0.080000000) (4.880000000,0.096800000) (4.960000000,0.115200000) (5.040000000,0.135200000) (5.120000000,0.156800000) (5.200000000,0.180000000) (5.280000000,0.204800000) (5.360000000,0.231200000) (5.440000000,0.259200000) (5.520000000,0.288800000) (5.600000000,0.320000000) (5.680000000,0.352800000) (5.760000000,0.387200000) (5.840000000,0.423200000) (5.920000000,0.460800000) (6.000000000,0.500000000) (6.080000000,0.540800000) (6.160000000,0.583200000) (6.240000000,0.627200000) (6.320000000,0.672800000) (6.400000000,0.720000000) (6.480000000,0.768800000) (6.560000000,0.819200000) (6.640000000,0.871200000) (6.720000000,0.924800000) (6.800000000,0.980000000) (6.880000000,1.036800000) (6.960000000,1.095200000) (7.040000000,1.155200000) (7.120000000,1.216800000) (7.200000000,1.280000000) (7.280000000,1.344800000) (7.360000000,1.411200000) (7.440000000,1.479200000) (7.520000000,1.548800000) (7.600000000,1.620000000) (7.680000000,1.692800000) (7.760000000,1.767200000) (7.840000000,1.843200000) (7.920000000,1.920800000) (8.000000000,2.000000000)};
\node[below] at (0,-.12) {$0$};
\node[below] at (4,-.12) {$1/2$};
\node[below] at (8,-.12) {$1$};
\draw[guide] (0,2)--(8.2,2) node[right] {$\Delta K=0$};\node[left] at (0,0) {$-1/2$};\node[left] at (0,2) {$0$};\end{tikzpicture}
\figcaption{解图力试2.4C\quad 累计冲量增大：总动能先降后回到初值}
\end{minipage}\end{center}
图C的横坐标用最终弹性冲量$J_{\rm el}$归一化，纵坐标表示$\mathcal K\Delta K/v_0^2$，水平虚线对应总动能回到碰前值。累计冲量从零增大时，刚体动能转入接触形变能，在$J=v_0/\mathcal K$处最低，此时两端法向速度相同；随后形变恢复，累计冲量继续增加，动能回到碰前值。由$\Delta K=0$，除去没有发生碰撞的$J=0$根，得
\[
\boxed{J=\frac{2v_0}{\mathcal K}
=\frac{m_Am_B}{2(m_A+m_B)}v_0}.
\]
于是
\[
\boxed{v_A=\frac{m_B}{2(m_A+m_B)}v_0,\qquad
v_B=\frac{m_A+2m_B}{2(m_A+m_B)}v_0},
\]
\[
\boxed{\omega_A=\frac{3m_B}{l(m_A+m_B)}v_0,\qquad
\omega_B=\frac{3m_A}{l(m_A+m_B)}v_0}.
\]
两杆质心都向上运动，两角速度也都为正。相反的接触冲量作用在相反侧的端部，因而绕各自质心都产生逆时针转动。

\begin{center}\begin{minipage}{.98\linewidth}\centering
\begin{tikzpicture}[scale=.83,transform shape]
\draw[very thick,color=accent] (-4.00000,0.00000)--(0.00000,0.00000);\draw[very thick,color=ink] (0.00000,-3.00000)--(4.00000,-3.00000);\draw[motion] (-2.00000,0.00000)--(-2.00000,0.60000);\node[above left] at (-2.00000,0.60000) {$v_A=v_0/4$};\draw[force] (0.00000,0.00000)--(0.00000,2.40000);\node[above right] at (0.00000,2.40000) {$u_A=v_0$};\draw[motion] (2.00000,-3.00000)--(2.00000,-1.20000);\node[above right] at (2.00000,-1.20000) {$v_B=3v_0/4$};\fill (0.00000,-3.00000) circle(.035);\node[below left] at (0.00000,-3.00000) {$u_B=0$};\node[above] at (0.00000,3.50000) {等质量示例：$m_A=m_B$};\node[below] at (0.00000,-4.40000) {恢复关系比较端点，不只比较质心};
\end{tikzpicture}
\figcaption{解图力试2.4B\quad 等质量时的质心与接触点速度，uB恰为零。}
\end{minipage}\end{center}
（2）解图力试2.4B把接触点速度与质心速度分开，箭头取等质量的情形举例。一般情况下，杆$A$右端和杆$B$左端的碰后法向速度分别为
\[
u_A=v_A+\frac l2\omega_A=\frac{4J}{m_A},
\qquad
u_B=v_B-\frac l2\omega_B=v_0-\frac{4J}{m_B}.
\]
所以
\[
u_A-u_B=\mathcal KJ-v_0=v_0,
\]
即
\[
\boxed{v_0=
\left(v_A+\frac l2\omega_A\right)
-\left(v_B-\frac l2\omega_B\right)}.
\]
碰前接近速率为$v_0$，碰后分离速率也为$v_0$，这是恢复系数为1的接触关系。这个关系由完整刚体的能量与冲量方程直接推出，也适用于这里的两杆端部碰撞。

本题取光滑单点法向接触，碰前、碰后均以完整刚体描述，完全弹性过程回收接触时暂存的形变能。端部质元受到杆内力，计算时把整杆的平动与转动一并计入，端点速度按$v\pm(l/2)\omega$合成。摩擦接触需增加切向冲量，多点接触需联立各接触约束，若振动仍储存能量，则把它计入碰后的能量分配。

总动量给$m_Av_A+m_Bv_B=m_Bv_0$。关于接触点的总角动量，则由两杆质心的轨道角动量与各自绕质心的自转角动量共同组成。
把这一总角动量展开，以接触点$O$为固定点，有
\[
\frac l2m_Bv_0=\frac l2m_Bv_B-\frac l2m_Av_A+I_A\omega_A+I_B\omega_B.
\]
式中$A$质心位于$O$左侧，故其向上平移的轨道项为负；$B$位于右侧，其轨道项为正。把冲量方程代入，右侧相对碰前的变化为$-Jl/2-Jl/2+Jl/2+Jl/2=0$，轨道角动量的减少恰好转为两杆的自转角动量。进一步由角冲量得$m_A\omega_A=m_B\omega_B$，与总动量、总角动量联立便得到$\omega_A=6v_A/l$、$\omega_B=6m_Av_A/(m_Bl)$和$v_B=v_0-m_Av_A/m_B$，与以$J$为变量的写法相同。

能量式的两根分别是$J=0$与$2v_0/\mathcal K$。碰前两接触端正在接近，光滑单侧接触会产生正法向冲量，所以弹性终态取第二根。此时$u_A-u_B=v_0>0$，两杆端部沿法向分开；质心速度差与端点分离速率分别对应不同的图中对象。接触时间趋零保证计算角冲量时可固定力臂，$\epsilon/l\to0$则把实际短重叠接触的位置化为所画的端部极限。



\Needspace{13\baselineskip}

\pdfbookmark[2]{第 5 题 对称铰接双杆剪绳后的下落}{q5}

\section*{第 5 题\quad 对称铰接双杆剪绳后的下落}\label{q:5}


两根相同均匀细杆各长$l$、质量各为$m$，上端以光滑铰链$C$连接，下端放在光滑水平地面上，成对称的人字形。起初每杆与水平面夹角$\theta_0=30^\circ$，用一根轻绳保持静止。突然剪断轻绳，双杆在竖直平面内运动，忽略阻力。
\begin{enumerate}
\item 求铰链$C$即将碰到地面时的速度。
\item 用积分式表示从剪绳到铰链碰地所需的时间，并检查两杆下端在这段时间内是否保持接地。
\end{enumerate}
只求第一次碰地之前的运动，不讨论碰地冲击。


\begin{center}\begin{minipage}{.98\linewidth}\centering
\begin{tikzpicture}[scale=.9,transform shape]
\draw[object] (-5.00000,0.00000)--(5.00000,0.00000);\draw[very thick,color=accent] (-3.46410,0.00000)--(0.00000,2.00000)--(3.46410,0.00000);\fill (0.00000,2.00000) circle(.035);\node[above] at (0.00000,2.00000) {$C$};\draw[guide] (-1.73205,1.00000)--(1.73205,1.00000);\node[above] at (0.00000,1.10000) {剪断轻绳};\node[above left] at (-2.60000,1.40000) {$l,m$};\node[above right] at (2.60000,1.40000) {$l,m$};\node[below] at (-2.90000,-0.10000) {$\theta_0$};\node[below] at (2.90000,-0.10000) {$\theta_0$};
\end{tikzpicture}
\figcaption{图力试2.5\quad 两杆对称，绳剪断后脚点向两侧张开，铰链竖直下降。}
\end{minipage}\end{center}

\Needspace{6\baselineskip}\subsection*{解答}


图力试2.5的左右两杆完全对称，铰链$C$的横坐标始终为零，只沿竖直方向运动。取地面为$y=0$，以铰链竖直投影为$x=0$。设当前杆与水平面的夹角为$\theta$，剪绳后$\theta$从$\theta_0$减小到0。

\begin{center}\begin{minipage}{.98\linewidth}\centering
\begin{tikzpicture}[scale=.83,transform shape]
\draw[axis] (0.00000,0.00000)--(5.25000,0.00000);\draw[axis] (0.00000,0.00000)--(0.00000,2.85000);\node[below right] at (5.25000,0.00000) {$x$};\node[above right] at (0.00000,2.85000) {$y$};\draw[very thick,color=accent] (0.00000,1.37159)--(3.75749,0.00000);\fill (0.00000,1.37159) circle(.035);\node[above left] at (0.00000,1.37159) {$C$};\fill (1.87875,0.68580) circle(.035);\node[above right] at (1.87875,0.68580) {$G$};\fill (3.75749,0.00000) circle(.035);\node[below] at (3.75749,0.00000) {脚点};\draw[guide] (0.00000,1.37159)--(-0.60000,1.37159);\draw[motion] (-0.60000,1.37159)--(-0.60000,0.20000);\node[left] at (-0.70000,0.60000) {$v_C$};\draw[motion] (1.87875,0.68580)--(2.09875,0.08580);\node[below left] at (1.80000,-0.10000) {$\boldsymbol v_G$};\draw[force] (3.75749,0.00000)--(3.75749,0.90000);\node[right] at (3.85749,0.80000) {$N$};\node[below] at (2.80000,-0.70000) {$\theta$};\draw[guide] (2.80000,-0.60000)--(3.30749,0.08000);\draw[guide] plot coordinates {(3.19387,0.20574) (3.19143,0.19892) (3.18907,0.19208) (3.18679,0.18520) (3.18460,0.17830) (3.18249,0.17137) (3.18046,0.16442) (3.17852,0.15745) (3.17666,0.15045) (3.17489,0.14343) (3.17320,0.13638) (3.17159,0.12932) (3.17008,0.12224) (3.16864,0.11514) (3.16730,0.10803) (3.16604,0.10090) (3.16486,0.09375) (3.16377,0.08659) (3.16277,0.07942) (3.16186,0.07224) (3.16103,0.06504) (3.16029,0.05784) (3.15963,0.05063) (3.15906,0.04341) (3.15858,0.03618) (3.15819,0.02895) (3.15788,0.02172) (3.15767,0.01448) (3.15753,0.00724) (3.15749,0.00000)};\node[below] at (2.00000,-1.40000) {$y_C=2y_G=l\sin\theta$};
\end{tikzpicture}
\figcaption{解图力试2.5A\quad 铰链、质心和脚点速度不同；脚点可在地面滑动。}
\end{minipage}\end{center}
解图力试2.5A只画右杆。其下端位于$(l\cos\theta,0)$，铰链位于$(0,l\sin\theta)$，杆质心坐标和速度为
\[
x_G=\frac l2\cos\theta,\quad y_G=\frac l2\sin\theta,
\]
\[
\dot x_G=-\frac l2\sin\theta\,\dot\theta,\qquad
\dot y_G=\frac l2\cos\theta\,\dot\theta.
\]
所以质心速率平方为$l^2\dot\theta^2/4$。加上杆绕自身质心的转动能，单杆动能为
\[
K_1=\frac12m\frac{l^2\dot\theta^2}{4}
+\frac12\frac{ml^2}{12}\dot\theta^2
=\frac16ml^2\dot\theta^2.
\]
两杆动能、势能都是单杆的两倍，因而可将总能量方程同时除以2：
\[
\frac16ml^2\dot\theta^2+\frac12mgl\sin\theta
=\frac12mgl\sin\theta_0.
\]
对能量式先求时间导数，得到$\dot\theta[ml^2\ddot\theta/3+mgl\cos\theta/2]=0$。沿开始下落后的连续运动支，$\ddot\theta=-3g\cos\theta/(2l)$；初始静止点的加速度取这个连续值。再用$\mathrm d(\dot\theta^2)=2\ddot\theta\,\mathrm d\theta$积分，得$\dot\theta^2=-(3g/l)\sin\theta+C_0$，初值$\dot\theta(0)=0$给$C_0=(3g/l)\sin\theta_0$。因此能量积分与角运动方程给出相同的速度：
\[
\dot\theta^2=\frac{3g}{l}(\sin\theta_0-\sin\theta),
\qquad
\dot\theta=-\sqrt{\frac{3g}{l}(\sin\theta_0-\sin\theta)}.
\]
剪绳后夹角减小，故取负根，双杆脚点向两侧滑开。

（1）铰链高度为$y_C=l\sin\theta$，速度为
\[
\boldsymbol v_C=l\cos\theta\,\dot\theta\,\hat y.
\]
在第一次接触地面之前$\theta\to0^+$，所以
\[
\boxed{\boldsymbol v_C=-\sqrt{3gl\sin\theta_0}\,\hat y
=-\sqrt{\frac32gl}\,\hat y}.
\]
速率为$\sqrt{3gl/2}$；负号只表示竖直向下。铰链高度是质心高度的两倍，对时间求导后，铰链竖直速度也是质心竖直速度的两倍。

\begin{center}\begin{minipage}{.98\linewidth}\centering
\begin{tikzpicture}[scale=.8,transform shape]
\draw[axis] (0,0)--(10.7,0) node[right] {$\theta$};
\draw[axis] (0,0)--(0,3.2) node[above] {$|\dot\theta|\sqrt{l/g},\ N/(mg)$};
\node[above] at (5,3.5) {实线为角速率，虚线为支持力；初态在右端};
\draw[traj] plot coordinates {(0.000000000,2.816913204) (0.083333333,2.804595208) (0.166666667,2.792223107) (0.250000000,2.779796417) (0.333333333,2.767314639) (0.416666667,2.754777264) (0.500000000,2.742183770) (0.583333333,2.729533622) (0.666666667,2.716826271) (0.750000000,2.704061154) (0.833333333,2.691237693) (0.916666667,2.678355298) (1.000000000,2.665413359) (1.083333333,2.652411254) (1.166666667,2.639348343) (1.250000000,2.626223970) (1.333333333,2.613037459) (1.416666667,2.599788118) (1.500000000,2.586475237) (1.583333333,2.573098083) (1.666666667,2.559655907) (1.750000000,2.546147936) (1.833333333,2.532573377) (1.916666667,2.518931413) (2.000000000,2.505221205) (2.083333333,2.491441891) (2.166666667,2.477592581) (2.250000000,2.463672361) (2.333333333,2.449680289) (2.416666667,2.435615397) (2.500000000,2.421476684) (2.583333333,2.407263122) (2.666666667,2.392973650) (2.750000000,2.378607174) (2.833333333,2.364162567) (2.916666667,2.349638666) (3.000000000,2.335034270) (3.083333333,2.320348139) (3.166666667,2.305578993) (3.250000000,2.290725513) (3.333333333,2.275786330) (3.416666667,2.260760034) (3.500000000,2.245645166) (3.583333333,2.230440214) (3.666666667,2.215143617) (3.750000000,2.199753757) (3.833333333,2.184268958) (3.916666667,2.168687485) (4.000000000,2.153007540) (4.083333333,2.137227256) (4.166666667,2.121344698) (4.250000000,2.105357857) (4.333333333,2.089264648) (4.416666667,2.073062902) (4.500000000,2.056750368) (4.583333333,2.040324702) (4.666666667,2.023783466) (4.750000000,2.007124122) (4.833333333,1.990344024) (4.916666667,1.973440418) (5.000000000,1.956410426) (5.083333333,1.939251048) (5.166666667,1.921959149) (5.250000000,1.904531454) (5.333333333,1.886964537) (5.416666667,1.869254811) (5.500000000,1.851398520) (5.583333333,1.833391727) (5.666666667,1.815230301) (5.750000000,1.796909905) (5.833333333,1.778425979) (5.916666667,1.759773728) (6.000000000,1.740948104) (6.083333333,1.721943783) (6.166666667,1.702755151) (6.250000000,1.683376276) (6.333333333,1.663800886) (6.416666667,1.644022342) (6.500000000,1.624033603) (6.583333333,1.603827199) (6.666666667,1.583395189) (6.750000000,1.562729121) (6.833333333,1.541819983) (6.916666667,1.520658156) (7.000000000,1.499233351) (7.083333333,1.477534545) (7.166666667,1.455549907) (7.250000000,1.433266712) (7.333333333,1.410671245) (7.416666667,1.387748696) (7.500000000,1.364483026) (7.583333333,1.340856830) (7.666666667,1.316851165) (7.750000000,1.292445356) (7.833333333,1.267616770) (7.916666667,1.242340546) (8.000000000,1.216589279) (8.083333333,1.190332646) (8.166666667,1.163536952) (8.250000000,1.136164585) (8.333333333,1.108173356) (8.416666667,1.079515677) (8.500000000,1.050137559) (8.583333333,1.019977336) (8.666666667,0.988964054) (8.750000000,0.957015405) (8.833333333,0.924035013) (8.916666667,0.889908838) (9.000000000,0.854500304) (9.083333333,0.817643545) (9.166666667,0.779133819) (9.250000000,0.738713451) (9.333333333,0.696050500) (9.416666667,0.650704911) (9.500000000,0.602071819) (9.583333333,0.549279744) (9.666666667,0.490990149) (9.750000000,0.424948111) (9.833333333,0.346753661) (9.916666667,0.245038949) (10.000000000,0.000000030)};
\draw[dashed,thick,color=warm] plot coordinates {(0.000000000,0.575000000) (0.083333333,0.567571816) (0.166666667,0.560340815) (0.250000000,0.553307119) (0.333333333,0.546470834) (0.416666667,0.539832050) (0.500000000,0.533390843) (0.583333333,0.527147274) (0.666666667,0.521101388) (0.750000000,0.515253217) (0.833333333,0.509602775) (0.916666667,0.504150065) (1.000000000,0.498895071) (1.083333333,0.493837765) (1.166666667,0.488978102) (1.250000000,0.484316023) (1.333333333,0.479851455) (1.416666667,0.475584308) (1.500000000,0.471514479) (1.583333333,0.467641848) (1.666666667,0.463966283) (1.750000000,0.460487634) (1.833333333,0.457205740) (1.916666667,0.454120421) (2.000000000,0.451231484) (2.083333333,0.448538722) (2.166666667,0.446041914) (2.250000000,0.443740820) (2.333333333,0.441635191) (2.416666667,0.439724759) (2.500000000,0.438009243) (2.583333333,0.436488348) (2.666666667,0.435161763) (2.750000000,0.434029163) (2.833333333,0.433090210) (2.916666667,0.432344549) (3.000000000,0.431791812) (3.083333333,0.431431616) (3.166666667,0.431263566) (3.250000000,0.431287248) (3.333333333,0.431502237) (3.416666667,0.431908094) (3.500000000,0.432504365) (3.583333333,0.433290581) (3.666666667,0.434266259) (3.750000000,0.435430904) (3.833333333,0.436784005) (3.916666667,0.438325037) (4.000000000,0.440053462) (4.083333333,0.441968727) (4.166666667,0.444070267) (4.250000000,0.446357501) (4.333333333,0.448829836) (4.416666667,0.451486665) (4.500000000,0.454327366) (4.583333333,0.457351305) (4.666666667,0.460557834) (4.750000000,0.463946291) (4.833333333,0.467516001) (4.916666667,0.471266276) (5.000000000,0.475196415) (5.083333333,0.479305702) (5.166666667,0.483593411) (5.250000000,0.488058799) (5.333333333,0.492701112) (5.416666667,0.497519585) (5.500000000,0.502513436) (5.583333333,0.507681873) (5.666666667,0.513024090) (5.750000000,0.518539270) (5.833333333,0.524226581) (5.916666667,0.530085180) (6.000000000,0.536114212) (6.083333333,0.542312807) (6.166666667,0.548680085) (6.250000000,0.555215154) (6.333333333,0.561917109) (6.416666667,0.568785031) (6.500000000,0.575817993) (6.583333333,0.583015052) (6.666666667,0.590375256) (6.750000000,0.597897641) (6.833333333,0.605581229) (6.916666667,0.613425032) (7.000000000,0.621428051) (7.083333333,0.629589275) (7.166666667,0.637907681) (7.250000000,0.646382235) (7.333333333,0.655011893) (7.416666667,0.663795598) (7.500000000,0.672732283) (7.583333333,0.681820870) (7.666666667,0.691060270) (7.750000000,0.700449383) (7.833333333,0.709987099) (7.916666667,0.719672298) (8.000000000,0.729503847) (8.083333333,0.739480604) (8.166666667,0.749601419) (8.250000000,0.759865128) (8.333333333,0.770270558) (8.416666667,0.780816529) (8.500000000,0.791501846) (8.583333333,0.802325309) (8.666666667,0.813285704) (8.750000000,0.824381812) (8.833333333,0.835612400) (8.916666667,0.846976228) (9.000000000,0.858472048) (9.083333333,0.870098599) (9.166666667,0.881854613) (9.250000000,0.893738814) (9.333333333,0.905749916) (9.416666667,0.917886625) (9.500000000,0.930147635) (9.583333333,0.942531637) (9.666666667,0.955037309) (9.750000000,0.967663322) (9.833333333,0.980408340) (9.916666667,0.993271017) (10.000000000,1.006250000)};
\node[below] at (0,-.12) {$0$};
\node[below] at (5,-.12) {$15^\circ$};
\node[below] at (10,-.12) {$30^\circ$};
\end{tikzpicture}
\figcaption{解图力试2.5B\quad 实线为角速率，虚线为支持力；初态在右端}
\end{minipage}\end{center}
（2）解图力试2.5B画出$\theta$减小过程中角速率的增长。由时间微元$\dd t=\dd\theta/\dot\theta$，
\[
\boxed{t_f=
\sqrt{\frac{l}{3g}}\int_0^{\pi/6}
\frac{\dd\theta}{\sqrt{\tfrac12-\sin\theta}}
=\sqrt{\frac{2l}{3g}}\int_0^{\pi/6}
\frac{\dd\theta}{\sqrt{1-2\sin\theta}}}.
\]
积分的初态端有平方根奇点，因为从静止起步。近$\theta_0$时
\[
\sin\theta_0-\sin\theta
\sim\cos\theta_0(\theta_0-\theta),
\]
故被积函数只按$(\theta_0-\theta)^{-1/2}$发散，从初态到任意较小夹角所需的时间仍有限。若数值计算，可令$\theta=\theta_0-u^2$消除端点奇性。

下面由地面对脚点的支持力确定接地状态。\begin{center}\begin{minipage}{.98\linewidth}\centering
\begin{tikzpicture}[scale=.8,transform shape]
\draw[guide] (-.8,0)--(5,0);\draw[very thick] (0,2.4)--(4.156922,0);\fill (0,2.4) circle(.04) node[above] {$C$};\fill (2.078461,1.2) circle(.04) node[above right] {$G$};\draw[force] (2.078461,1.2)--(2.078461,-.2) node[below] {$mg$};\draw[force] (4.156922,0)--(4.156922,1.5) node[above] {$N$};\draw[force] (0,2.4)--(1.4,2.4) node[right] {$H$正向};\node at (2,-1.5) {$N-mg=m\ddot y_G,\quad H=m\ddot x_G$};
\end{tikzpicture}
\figcaption{解图力试2.5C\quad 右杆受力；铰链竖直分量由对称性为零}
\end{minipage}\end{center}
镜像对称使铰链对左右杆的竖直力相等，内力的作用反作用又使它们互为相反数，故每一竖直分量都是零。水平分量互为相反数，随运动推动或阻碍脚点向外滑动。图C的铰链速度竖直、力水平，它们垂直；脚点的支持力竖直、速度水平，也相互垂直，所以逐杆的能量守恒同样成立。对单杆竖直方向，
\[
N-mg=m\ddot y_G.
\]
由能量式的时间导数得$\ddot\theta=-3g\cos\theta/(2l)$，再对$y_G$求导，得到
\[
\boxed{N=\frac{mg}{4}
\left[1-6\sin\theta\sin\theta_0+9\sin^2\theta\right]}.
\]
这里先算$\ddot y_G=(l/2)(\cos\theta\ddot\theta-\sin\theta\dot\theta^2)$，再代入角加速度和能量积分，得到$\ddot y_G/g=-3/4-(3/2)\sin\theta\sin\theta_0+(9/4)\sin^2\theta$。这同时给出支持力曲线中的常数项、一次项和二次项。将括号配方：
\[
N=\frac{mg}{4}
\left[9\left(\sin\theta-\frac{\sin\theta_0}{3}\right)^2
+\cos^2\theta_0\right]
\ge\frac{mg}{4}\cos^2\theta_0.
\]
本题最小值为$3mg/16>0$，全程下端都保持接地，上述对称单自由度模型成立。铰链首次碰地后会受到新的冲量，随后的运动按碰地后的速度与接触状态计算。
本题初始支持力为$7mg/16$，在$\sin\theta=1/6$时降至最小$3mg/16$，随后回升，在铰链接地前趋向$mg/4$。这些值与图B的虚线对应，横轴右端表示刚剪绳，左端表示首次接地前。正支持力说明每个脚点始终受到地面的向上推力，双杆的人字形几何保持到铰链接地。角速率实线在右端从零起步，在左端达到$\sqrt{3g/(2l)}$；铰链速度还乘随角度变化的$l\cos\theta$，其方向始终向下。



\Needspace{13\baselineskip}

\pdfbookmark[2]{第 6 题 双线悬圆柱的平移、复摆与偏航}{q6}

\section*{第 6 题\quad 双线悬圆柱的平移、复摆与偏航}\label{q:6}


均匀实心圆柱质量为$m$、半径为$R$、轴长为$L$，由两根等长为$l$的轻细线对称悬挂。两线沿柱轴方向相距$D$，$0<D\le L$，下端固定在圆柱上表面；上端固定点为$O,O'$。圆柱质心$C$位于两线中间正下方。

求下列小振幅运动的周期，并明确各自的约束：
\begin{enumerate}
\item 两线在包含柱轴的同一铅垂平面内同角摆动，圆柱保持平行平移。
\item 把圆柱与悬线作为绕$OO'$轴摆动的一自由度受限复摆，求周期；并判断仅用柔线自由悬挂时，这种整体绕轴运动是否自动成立。
\item 圆柱在水平面内绕通过质心的铅垂轴小角偏航，求周期，并由薄圆片积分求所需转动惯量。
\end{enumerate}
忽略线的质量、伸长和阻力。受限复摆模式的导向不作功，其余运动由张紧柔线实现。


\begin{center}\begin{minipage}{.98\linewidth}\centering
\begin{tikzpicture}[scale=.9,transform shape]
\draw[object] plot coordinates {(-2.78000,-3.00000) (-2.78044,-2.96195) (-2.78177,-2.92404) (-2.78398,-2.88645) (-2.78705,-2.84931) (-2.79098,-2.81278) (-2.79576,-2.77700) (-2.80136,-2.74212) (-2.80775,-2.70828) (-2.81492,-2.67562) (-2.82284,-2.64426) (-2.83147,-2.61433) (-2.84078,-2.58595) (-2.85073,-2.55924) (-2.86128,-2.53431) (-2.87239,-2.51125) (-2.88401,-2.49016) (-2.89610,-2.47113) (-2.90861,-2.45422) (-2.92149,-2.43951) (-2.93468,-2.42706) (-2.94813,-2.41691) (-2.96180,-2.40912) (-2.97562,-2.40370) (-2.98953,-2.40068) (-3.00349,-2.40008) (-3.01743,-2.40189) (-3.03131,-2.40611) (-3.04506,-2.41272) (-3.05862,-2.42169) (-3.07195,-2.43300) (-3.08500,-2.44659) (-3.09769,-2.46240) (-3.11000,-2.48038) (-3.12186,-2.50046) (-3.13323,-2.52254) (-3.14407,-2.54655) (-3.15432,-2.57238) (-3.16396,-2.59994) (-3.17293,-2.62910) (-3.18121,-2.65976) (-3.18876,-2.69179) (-3.19554,-2.72506) (-3.20154,-2.75944) (-3.20673,-2.79479) (-3.21109,-2.83096) (-3.21459,-2.86781) (-3.21724,-2.90520) (-3.21900,-2.94297) (-3.21989,-2.98096) (-3.21989,-3.01904) (-3.21900,-3.05703) (-3.21724,-3.09480) (-3.21459,-3.13219) (-3.21109,-3.16904) (-3.20673,-3.20521) (-3.20154,-3.24056) (-3.19554,-3.27494) (-3.18876,-3.30821) (-3.18121,-3.34024) (-3.17293,-3.37090) (-3.16396,-3.40006) (-3.15432,-3.42762) (-3.14407,-3.45345) (-3.13323,-3.47746) (-3.12186,-3.49954) (-3.11000,-3.51962) (-3.09769,-3.53760) (-3.08500,-3.55341) (-3.07195,-3.56700) (-3.05862,-3.57831) (-3.04506,-3.58728) (-3.03131,-3.59389) (-3.01743,-3.59811) (-3.00349,-3.59992) (-2.98953,-3.59932) (-2.97562,-3.59630) (-2.96180,-3.59088) (-2.94813,-3.58309) (-2.93468,-3.57294) (-2.92149,-3.56049) (-2.90861,-3.54578) (-2.89610,-3.52887) (-2.88401,-3.50984) (-2.87239,-3.48875) (-2.86128,-3.46569) (-2.85073,-3.44076) (-2.84078,-3.41405) (-2.83147,-3.38567) (-2.82284,-3.35574) (-2.81492,-3.32438) (-2.80775,-3.29172) (-2.80136,-3.25788) (-2.79576,-3.22300) (-2.79098,-3.18722) (-2.78705,-3.15069) (-2.78398,-3.11355) (-2.78177,-3.07596) (-2.78044,-3.03805) (-2.78000,-3.00000)};\draw[object] plot coordinates {(3.22000,-3.00000) (3.21956,-2.96195) (3.21823,-2.92404) (3.21602,-2.88645) (3.21295,-2.84931) (3.20902,-2.81278) (3.20424,-2.77700) (3.19864,-2.74212) (3.19225,-2.70828) (3.18508,-2.67562) (3.17716,-2.64426) (3.16853,-2.61433) (3.15922,-2.58595) (3.14927,-2.55924) (3.13872,-2.53431) (3.12761,-2.51125) (3.11599,-2.49016) (3.10390,-2.47113) (3.09139,-2.45422) (3.07851,-2.43951) (3.06532,-2.42706) (3.05187,-2.41691) (3.03820,-2.40912) (3.02438,-2.40370) (3.01047,-2.40068) (2.99651,-2.40008) (2.98257,-2.40189) (2.96869,-2.40611) (2.95494,-2.41272) (2.94138,-2.42169) (2.92805,-2.43300) (2.91500,-2.44659) (2.90231,-2.46240) (2.89000,-2.48038) (2.87814,-2.50046) (2.86677,-2.52254) (2.85593,-2.54655) (2.84568,-2.57238) (2.83604,-2.59994) (2.82707,-2.62910) (2.81879,-2.65976) (2.81124,-2.69179) (2.80446,-2.72506) (2.79846,-2.75944) (2.79327,-2.79479) (2.78891,-2.83096) (2.78541,-2.86781) (2.78276,-2.90520) (2.78100,-2.94297) (2.78011,-2.98096) (2.78011,-3.01904) (2.78100,-3.05703) (2.78276,-3.09480) (2.78541,-3.13219) (2.78891,-3.16904) (2.79327,-3.20521) (2.79846,-3.24056) (2.80446,-3.27494) (2.81124,-3.30821) (2.81879,-3.34024) (2.82707,-3.37090) (2.83604,-3.40006) (2.84568,-3.42762) (2.85593,-3.45345) (2.86677,-3.47746) (2.87814,-3.49954) (2.89000,-3.51962) (2.90231,-3.53760) (2.91500,-3.55341) (2.92805,-3.56700) (2.94138,-3.57831) (2.95494,-3.58728) (2.96869,-3.59389) (2.98257,-3.59811) (2.99651,-3.59992) (3.01047,-3.59932) (3.02438,-3.59630) (3.03820,-3.59088) (3.05187,-3.58309) (3.06532,-3.57294) (3.07851,-3.56049) (3.09139,-3.54578) (3.10390,-3.52887) (3.11599,-3.50984) (3.12761,-3.48875) (3.13872,-3.46569) (3.14927,-3.44076) (3.15922,-3.41405) (3.16853,-3.38567) (3.17716,-3.35574) (3.18508,-3.32438) (3.19225,-3.29172) (3.19864,-3.25788) (3.20424,-3.22300) (3.20902,-3.18722) (3.21295,-3.15069) (3.21602,-3.11355) (3.21823,-3.07596) (3.21956,-3.03805) (3.22000,-3.00000)};\draw[object] (-3.00000,-2.40000)--(3.00000,-2.40000);\draw[object] (-3.00000,-3.60000)--(3.00000,-3.60000);\draw[object] (-2.30000,0.00000)--(2.30000,0.00000);\fill (-1.70000,0.00000) circle(.035);\node[above] at (-1.70000,0.00000) {$O$};\draw[object] (-1.70000,0.00000)--(-1.70000,-2.40000);\fill (1.70000,0.00000) circle(.035);\node[above] at (1.70000,0.00000) {$O'$};\draw[object] (1.70000,0.00000)--(1.70000,-2.40000);\fill (0.00000,-3.00000) circle(.035);\node[below] at (0.00000,-3.00000) {$C$};\draw[<->,thin,color=muted] (-1.70000,0.90000)--(1.70000,0.90000);\node[below] at (0.00000,0.90000) {$D$};\draw[<->,thin,color=muted] (4.00000,0.00000)--(4.00000,-2.40000);\node[right] at (4.00000,-1.20000) {$l$};\draw[<->,thin,color=muted] (-4.00000,-3.00000)--(-4.00000,-2.40000);\node[left] at (-4.00000,-2.70000) {$R$};\draw[<->,thin,color=muted] (-3.00000,-4.50000)--(3.00000,-4.50000);\node[below] at (0.00000,-4.50000) {$L$};
\end{tikzpicture}
\figcaption{图力试2.6\quad 柱轴方向两吊点相距D，线长l，质心低于吊点R。}
\end{minipage}\end{center}

\Needspace{6\baselineskip}\subsection*{解答}


图力试2.6中，线长$l$与质心到上悬轴的距离$l+R$不同。分别计算三类运动的动能和质心上升量，再据各自的约束确定周期。

\begin{center}\begin{minipage}{.98\linewidth}\centering
\begin{tikzpicture}[scale=.77,transform shape]
\draw[object] plot coordinates {(-2.12531,-3.52769) (-2.12575,-3.48964) (-2.12708,-3.45174) (-2.12929,-3.41414) (-2.13236,-3.37700) (-2.13630,-3.34047) (-2.14107,-3.30469) (-2.14667,-3.26982) (-2.15306,-3.23597) (-2.16024,-3.20331) (-2.16815,-3.17195) (-2.17678,-3.14202) (-2.18609,-3.11364) (-2.19604,-3.08694) (-2.20659,-3.06200) (-2.21770,-3.03895) (-2.22932,-3.01786) (-2.24141,-2.99882) (-2.25392,-2.98191) (-2.26680,-2.96720) (-2.27999,-2.95475) (-2.29344,-2.94461) (-2.30711,-2.93681) (-2.32093,-2.93139) (-2.33484,-2.92837) (-2.34880,-2.92777) (-2.36275,-2.92958) (-2.37662,-2.93380) (-2.39037,-2.94041) (-2.40394,-2.94939) (-2.41727,-2.96069) (-2.43031,-2.97428) (-2.44301,-2.99010) (-2.45531,-3.00808) (-2.46717,-3.02815) (-2.47855,-3.05024) (-2.48938,-3.07424) (-2.49964,-3.10008) (-2.50927,-3.12763) (-2.51824,-3.15680) (-2.52652,-3.18746) (-2.53407,-3.21949) (-2.54085,-3.25276) (-2.54685,-3.28713) (-2.55204,-3.32248) (-2.55640,-3.35865) (-2.55991,-3.39551) (-2.56255,-3.43289) (-2.56431,-3.47066) (-2.56520,-3.50866) (-2.56520,-3.54673) (-2.56431,-3.58473) (-2.56255,-3.62249) (-2.55991,-3.65988) (-2.55640,-3.69673) (-2.55204,-3.73290) (-2.54685,-3.76825) (-2.54085,-3.80263) (-2.53407,-3.83590) (-2.52652,-3.86793) (-2.51824,-3.89859) (-2.50927,-3.92775) (-2.49964,-3.95531) (-2.48938,-3.98114) (-2.47855,-4.00515) (-2.46717,-4.02723) (-2.45531,-4.04731) (-2.44301,-4.06529) (-2.43031,-4.08110) (-2.41727,-4.09469) (-2.40394,-4.10600) (-2.39037,-4.11497) (-2.37662,-4.12159) (-2.36275,-4.12581) (-2.34880,-4.12762) (-2.33484,-4.12701) (-2.32093,-4.12400) (-2.30711,-4.11858) (-2.29344,-4.11078) (-2.27999,-4.10063) (-2.26680,-4.08818) (-2.25392,-4.07347) (-2.24141,-4.05656) (-2.22932,-4.03753) (-2.21770,-4.01644) (-2.20659,-3.99338) (-2.19604,-3.96845) (-2.18609,-3.94174) (-2.17678,-3.91336) (-2.16815,-3.88344) (-2.16024,-3.85208) (-2.15306,-3.81941) (-2.14667,-3.78557) (-2.14107,-3.75069) (-2.13630,-3.71491) (-2.13236,-3.67838) (-2.12929,-3.64124) (-2.12708,-3.60365) (-2.12575,-3.56575) (-2.12531,-3.52769)};\draw[object] plot coordinates {(3.87469,-3.52769) (3.87425,-3.48964) (3.87292,-3.45174) (3.87071,-3.41414) (3.86764,-3.37700) (3.86370,-3.34047) (3.85893,-3.30469) (3.85333,-3.26982) (3.84694,-3.23597) (3.83976,-3.20331) (3.83185,-3.17195) (3.82322,-3.14202) (3.81391,-3.11364) (3.80396,-3.08694) (3.79341,-3.06200) (3.78230,-3.03895) (3.77068,-3.01786) (3.75859,-2.99882) (3.74608,-2.98191) (3.73320,-2.96720) (3.72001,-2.95475) (3.70656,-2.94461) (3.69289,-2.93681) (3.67907,-2.93139) (3.66516,-2.92837) (3.65120,-2.92777) (3.63725,-2.92958) (3.62338,-2.93380) (3.60963,-2.94041) (3.59606,-2.94939) (3.58273,-2.96069) (3.56969,-2.97428) (3.55699,-2.99010) (3.54469,-3.00808) (3.53283,-3.02815) (3.52145,-3.05024) (3.51062,-3.07424) (3.50036,-3.10008) (3.49073,-3.12763) (3.48176,-3.15680) (3.47348,-3.18746) (3.46593,-3.21949) (3.45915,-3.25276) (3.45315,-3.28713) (3.44796,-3.32248) (3.44360,-3.35865) (3.44009,-3.39551) (3.43745,-3.43289) (3.43569,-3.47066) (3.43480,-3.50866) (3.43480,-3.54673) (3.43569,-3.58473) (3.43745,-3.62249) (3.44009,-3.65988) (3.44360,-3.69673) (3.44796,-3.73290) (3.45315,-3.76825) (3.45915,-3.80263) (3.46593,-3.83590) (3.47348,-3.86793) (3.48176,-3.89859) (3.49073,-3.92775) (3.50036,-3.95531) (3.51062,-3.98114) (3.52145,-4.00515) (3.53283,-4.02723) (3.54469,-4.04731) (3.55699,-4.06529) (3.56969,-4.08110) (3.58273,-4.09469) (3.59606,-4.10600) (3.60963,-4.11497) (3.62338,-4.12159) (3.63725,-4.12581) (3.65120,-4.12762) (3.66516,-4.12701) (3.67907,-4.12400) (3.69289,-4.11858) (3.70656,-4.11078) (3.72001,-4.10063) (3.73320,-4.08818) (3.74608,-4.07347) (3.75859,-4.05656) (3.77068,-4.03753) (3.78230,-4.01644) (3.79341,-3.99338) (3.80396,-3.96845) (3.81391,-3.94174) (3.82322,-3.91336) (3.83185,-3.88344) (3.83976,-3.85208) (3.84694,-3.81941) (3.85333,-3.78557) (3.85893,-3.75069) (3.86370,-3.71491) (3.86764,-3.67838) (3.87071,-3.64124) (3.87292,-3.60365) (3.87425,-3.56575) (3.87469,-3.52769)};\draw[object] (-2.34531,-2.92769)--(3.65469,-2.92769);\draw[object] (-2.34531,-4.12769)--(3.65469,-4.12769);\draw[object] (-2.50000,0.00000)--(2.50000,0.00000);\draw[object] (-2.00000,0.00000)--(-1.34531,-2.92769);\draw[guide] (-2.00000,0.00000)--(-2.00000,-3.00000);\draw[object] (2.00000,0.00000)--(2.65469,-2.92769);\draw[guide] (2.00000,0.00000)--(2.00000,-3.00000);\fill (0.65469,-3.52769) circle(.035);\node[below] at (0.65469,-3.52769) {$C$};\draw[motion] (0.65469,-3.52769)--(1.85469,-3.25935);\node[right] at (1.95469,-3.22769) {$l\dot q$};\node[above] at (-1.50000,-1.00000) {$q$};\node[above] at (0.00000,1.10000) {柱轴保持平行，圆柱不转};\node[below] at (0.00000,-5.30000) {$\Delta h=l(1-\cos q)$};
\end{tikzpicture}
\figcaption{解图力试2.6A\quad 同角平移时，质心上升由线长l决定。}
\end{minipage}\end{center}
（1）解图力试2.6A表示平移模式：两线保持平行，均偏转小角$q$，圆柱姿态不变。质心与下吊点有相同的平移速度$l\dot q$，质心相对初态升高$l(1-\cos q)$，于是
\[
K=\frac12ml^2\dot q^2,\qquad
U=mgl(1-\cos q).
\]
小角下$m l^2\ddot q+mglq=0$，所以
\[
\boxed{T_1=2\pi\sqrt{\frac lg}}.
\]
两上悬点的间距与两下吊点的间距同为$D$，两线等长且同向摆动，沿柱轴的平面内构成平行四边形，所以两下吊点的连接线保持水平。以$q$对能量求导，得到$ml^2\ddot q+mgl\sin q=0$；在线性近似中用$\sin q\simeq q$得到上述周期。这里圆柱不转，质心虽比线下端低$R$，其平移圆弧半径仍为$l$。两线拉力的差调节绕质心的俯仰力矩，使圆柱保持姿态。

\begin{center}\begin{minipage}{.98\linewidth}\centering
\begin{tikzpicture}[scale=.8,transform shape]
\draw[object] (1.37159,-3.75749) circle(1.00000);\draw[object] (0.00000,0.00000)--(1.02869,-2.81812);\draw[guide] (1.02869,-2.81812)--(1.37159,-3.75749);\draw[guide] (0.00000,0.00000)--(0.00000,-4.80000);\fill (0.00000,0.00000) circle(.035);\node[above] at (0.00000,0.00000) {$OO'$悬轴};\fill (1.02869,-2.81812) circle(.035);\draw[guide] (1.02869,-2.81812)--(2.40000,-2.60000);\node[right] at (2.50000,-2.60000) {吊点};\fill (1.37159,-3.75749) circle(.035);\node[below left] at (1.37159,-3.75749) {$C$};\node[above] at (0.50000,-0.80000) {$\chi$};\node[right] at (2.30000,-1.80000) {$d=l+R$};\node[above] at (1.00000,1.20000) {受限模式：$q=\psi=\chi$};\node[below] at (1.00000,-5.80000) {理想导向保持两角相同};
\end{tikzpicture}
\figcaption{解图力试2.6B\quad 规定受限复摆用上悬轴惯量和距离l+R。}
\end{minipage}\end{center}
（2）解图力试2.6B先按规定的受限模式，令线与圆柱作为整体绕固定轴$OO'$摆动。质心到轴距离为$d=l+R$，圆柱绕自身柱轴的惯量为$I_\parallel=mR^2/2$，平行轴定理给
\[
I_{OO'}=\frac12mR^2+md^2.
\]
以整体摆角$\chi$为坐标，
\[
K=\frac12I_{OO'}\dot\chi^2,\qquad
U=mgd(1-\cos\chi).
\]
小角恢复力矩为$-mgd\chi$，故
\[
\boxed{T_2=2\pi
\sqrt{\frac{\tfrac12R^2+(l+R)^2}{g(l+R)}}}.
\]
此式对应的是线方向与圆柱滚转方向相同的一自由度受限运动。同步转动由理想导向维持；当柔线拉力与质心、上悬轴共线时，拉力对圆柱质心的滚转力矩为零，而同步复摆需要非零滚转角加速度，因此导向还提供相应的约束力矩。

对于柔线自由悬挂，解图力试2.6C分开横向线摆角$q$和圆柱绕柱轴滚转角$\psi$。在垂直柱轴的截面中，质心坐标可取
\[
y_C=l\sin q+R\sin\psi,\qquad
z_C=-l\cos q-R\cos\psi.
\]
\begin{center}\begin{minipage}{.98\linewidth}\centering
\begin{tikzpicture}[scale=.8,transform shape]
\draw[object] (0.54366,-3.80538) circle(1.00000);\draw[object] (0.00000,0.00000)--(0.88656,-2.86601);\draw[traj] (0.88656,-2.86601)--(0.54366,-3.80538);\draw[guide] (0.00000,0.00000)--(0.00000,-5.00000);\draw[guide] (0.88656,-2.86601)--(0.88656,-4.46601);\fill (0.00000,0.00000) circle(.035);\node[above] at (0.00000,0.00000) {上悬轴};\fill (0.88656,-2.86601) circle(.035);\node[right] at (0.88656,-2.86601) {吊点};\fill (0.54366,-3.80538) circle(.035);\node[below left] at (0.54366,-3.80538) {$C$};\node[above] at (0.45000,-0.80000) {$q$};\node[left] at (0.63656,-3.46601) {$\psi$};\draw[force] (0.88656,-2.86601)--(0.44328,-1.43300);\node[left] at (-0.40000,-1.70000) {$F$};\draw[guide] (-0.30000,-1.70000)--(0.55000,-1.80000);\node[above] at (1.00000,1.20000) {柔线：线角与本体滚转角独立};\node[below] at (1.00000,-6.10000) {拉力通常不通过C，产生滚转力矩};
\end{tikzpicture}
\figcaption{解图力试2.6C\quad 柔线横摆与圆柱滚转独立，须求两个耦合模态。}
\end{minipage}\end{center}
小角二阶动能、势能为
\[
K=\frac m2\left(l^2\dot q^2+2lR\dot q\dot\psi
+\frac32R^2\dot\psi^2\right),\qquad
U=\frac{mg}{2}(lq^2+R\psi^2).
\]
质心速度由两段运动相加，$\dot y_C\simeq l\dot q+R\dot\psi$；竖直速度在平衡点附近为更高阶。质心平动能是$m(l\dot q+R\dot\psi)^2/2$，再加柱轴转动能$mR^2\dot\psi^2/4$，得到$\dot\psi^2$项系数$3mR^2/4$。将$U$分别对$q,\psi$求导，势能没有交叉项，动能却含交叉项。因此自由横摆与滚转由两个坐标耦合描述：
\[
\begin{pmatrix}l^2&lR\\lR&3R^2/2\end{pmatrix}
\begin{pmatrix}\ddot q\\\ddot\psi\end{pmatrix}
+g\begin{pmatrix}l&0\\0&R\end{pmatrix}
\begin{pmatrix}q\\\psi\end{pmatrix}=0.
\]
逐个写出的方程是$l^2\ddot q+lR\ddot\psi+glq=0$与$lR\ddot q+(3/2)R^2\ddot\psi+gR\psi=0$。令解按$\exp(i\omega t)$变化，系数行列式为$(gl-l^2\omega^2)(gR-3R^2\omega^2/2)-(lR\omega^2)^2=0$；除以正数$lR$，行列式给
\[
\frac12lR\,\omega^4
-g\left(l+\frac32R\right)\omega^2+g^2=0,
\]
\[
\boxed{\omega_\pm^2=
\frac{g\left[2l+3R\pm\sqrt{4l^2+4lR+9R^2}\right]}{2lR}},
\qquad T_\pm=\frac{2\pi}{\omega_\pm}.
\]
两个自由模态各有自己的线摆与滚转比例。附加理想导向若强制$q=\psi=\chi$，两角速度也相同，动能变为$m[(l+R)^2+R^2/2]\dot\chi^2/2$，势能变为$mg(l+R)\chi^2/2$，由二次能量之比恢复上面的受限周期。自由模态的角度比例则由质量矩阵第一行给出$\psi/q=g/(R\omega^2)-l/R$。

\begin{center}\begin{minipage}{.98\linewidth}\centering
\begin{tikzpicture}[scale=.68,transform shape]
\draw[guide] (-3.00000,-0.50000)--(3.00000,-0.50000)--(3.00000,0.50000)--(-3.00000,0.50000)--(-3.00000,-0.50000);\draw[object] (-2.71825,-1.36423)--(3.01377,0.40889)--(2.71825,1.36423)--(-3.01377,-0.40889)--(-2.71825,-1.36423);\fill (0.00000,0.00000) circle(.035);\node[below] at (0.00000,0.00000) {$C$};\fill (-2.00000,0.00000) circle(.035);\fill (-1.91067,-0.59104) circle(.035);\draw[force] (-1.91067,-0.59104)--(-2.00000,0.00000);\draw[guide] (-2.00000,0.00000)--(-2.00000,0.90000);\node[above left] at (-2.00000,0.90000) {$O$};\fill (2.00000,0.00000) circle(.035);\fill (1.91067,0.59104) circle(.035);\draw[force] (1.91067,0.59104)--(2.00000,0.00000);\draw[guide] (2.00000,0.00000)--(2.00000,-1.00000);\node[below right] at (2.00000,-1.00000) {$O'$};\draw[guide] plot coordinates {(1.00000,0.00000) (0.99995,0.01034) (0.99979,0.02069) (0.99952,0.03103) (0.99914,0.04137) (0.99866,0.05170) (0.99807,0.06203) (0.99738,0.07235) (0.99658,0.08266) (0.99567,0.09297) (0.99465,0.10326) (0.99353,0.11355) (0.99230,0.12382) (0.99097,0.13408) (0.98953,0.14432) (0.98798,0.15455) (0.98633,0.16476) (0.98458,0.17496) (0.98271,0.18513) (0.98075,0.19529) (0.97867,0.20542) (0.97650,0.21554) (0.97421,0.22563) (0.97183,0.23569) (0.96934,0.24573) (0.96674,0.25575) (0.96405,0.26573) (0.96125,0.27569) (0.95834,0.28562) (0.95534,0.29552)};\node[below] at (1.10000,-1.30000) {$\theta$};\draw[guide] (1.10000,-1.20000)--(0.95000,0.15000);\node[right] at (3.80000,1.30000) {侧移 $\simeq D\theta/2$};\draw[guide] (3.70000,1.20000)--(1.96000,0.30000);\begin{scope}[shift={(9.00000,0.00000)}]
\draw[guide] (0.00000,0.00000)--(0.00000,-3.00000);\draw[object] (0.00000,0.00000)--(0.80000,-2.50000);\fill (0.80000,-2.50000) circle(.035);\node[right] at (0.80000,-2.50000) {下吊点};\node[right] at (0.30000,-0.80000) {$\phi$};\node[above] at (0.60000,1.00000) {线的侧视};\node[below] at (0.50000,-4.00000) {$\phi\simeq D\theta/(2l)$};
\end{scope}\node[above] at (0.00000,2.00000) {偏航俯视；张力水平分量指向固定点};
\end{tikzpicture}
\figcaption{解图力试2.6D\quad 线侧移产生二阶上升量，决定线性偏航恢复力。}
\end{minipage}\end{center}
（3）解图力试2.6D表示水平偏航。圆柱转过小角$\theta$时，每个下吊点相对正上方固定点的水平位移约为$D\theta/2$，故线偏角为
\[
\phi\simeq\frac{D}{2l}\theta.
\]
两线等长使质心升高
\[
\Delta h=l(1-\cos\phi)
\simeq\frac{D^2}{8l}\theta^2.
\]
俯视图中每个下吊点到其正上方的位移是小弧长$D\theta/2$，线长保持$l$，下吊点与质心同时升高。于是高度近似式中的因子$1/8$来自$1-\cos\phi\simeq\phi^2/2$与力臂$D/2$的平方。所以势能与恢复系数为
\[
U\simeq\frac{mgD^2}{8l}\theta^2
=\frac12k_\theta\theta^2,\qquad
k_\theta=\frac{mgD^2}{4l}.
\]
\begin{center}\begin{minipage}{.98\linewidth}\centering
\begin{tikzpicture}[scale=.73,transform shape]
\draw[object] plot coordinates {(-2.78000,0.00000) (-2.78044,0.05074) (-2.78177,0.10127) (-2.78398,0.15140) (-2.78705,0.20092) (-2.79098,0.24963) (-2.79576,0.29733) (-2.80136,0.34384) (-2.80775,0.38896) (-2.81492,0.43251) (-2.82284,0.47433) (-2.83147,0.51423) (-2.84078,0.55206) (-2.85073,0.58767) (-2.86128,0.62092) (-2.87239,0.65166) (-2.88401,0.67978) (-2.89610,0.70516) (-2.90861,0.72771) (-2.92149,0.74732) (-2.93468,0.76392) (-2.94813,0.77745) (-2.96180,0.78785) (-2.97562,0.79507) (-2.98953,0.79909) (-3.00349,0.79990) (-3.01743,0.79748) (-3.03131,0.79186) (-3.04506,0.78304) (-3.05862,0.77107) (-3.07195,0.75600) (-3.08500,0.73788) (-3.09769,0.71680) (-3.11000,0.69282) (-3.12186,0.66606) (-3.13323,0.63661) (-3.14407,0.60460) (-3.15432,0.57016) (-3.16396,0.53342) (-3.17293,0.49453) (-3.18121,0.45365) (-3.18876,0.41094) (-3.19554,0.36658) (-3.20154,0.32074) (-3.20673,0.27362) (-3.21109,0.22539) (-3.21459,0.17625) (-3.21724,0.12640) (-3.21900,0.07604) (-3.21989,0.02538) (-3.21989,-0.02538) (-3.21900,-0.07604) (-3.21724,-0.12640) (-3.21459,-0.17625) (-3.21109,-0.22539) (-3.20673,-0.27362) (-3.20154,-0.32074) (-3.19554,-0.36658) (-3.18876,-0.41094) (-3.18121,-0.45365) (-3.17293,-0.49453) (-3.16396,-0.53342) (-3.15432,-0.57016) (-3.14407,-0.60460) (-3.13323,-0.63661) (-3.12186,-0.66606) (-3.11000,-0.69282) (-3.09769,-0.71680) (-3.08500,-0.73788) (-3.07195,-0.75600) (-3.05862,-0.77107) (-3.04506,-0.78304) (-3.03131,-0.79186) (-3.01743,-0.79748) (-3.00349,-0.79990) (-2.98953,-0.79909) (-2.97562,-0.79507) (-2.96180,-0.78785) (-2.94813,-0.77745) (-2.93468,-0.76392) (-2.92149,-0.74732) (-2.90861,-0.72771) (-2.89610,-0.70516) (-2.88401,-0.67978) (-2.87239,-0.65166) (-2.86128,-0.62092) (-2.85073,-0.58767) (-2.84078,-0.55206) (-2.83147,-0.51423) (-2.82284,-0.47433) (-2.81492,-0.43251) (-2.80775,-0.38896) (-2.80136,-0.34384) (-2.79576,-0.29733) (-2.79098,-0.24963) (-2.78705,-0.20092) (-2.78398,-0.15140) (-2.78177,-0.10127) (-2.78044,-0.05074) (-2.78000,-0.00000)};\draw[object] plot coordinates {(3.22000,0.00000) (3.21956,0.05074) (3.21823,0.10127) (3.21602,0.15140) (3.21295,0.20092) (3.20902,0.24963) (3.20424,0.29733) (3.19864,0.34384) (3.19225,0.38896) (3.18508,0.43251) (3.17716,0.47433) (3.16853,0.51423) (3.15922,0.55206) (3.14927,0.58767) (3.13872,0.62092) (3.12761,0.65166) (3.11599,0.67978) (3.10390,0.70516) (3.09139,0.72771) (3.07851,0.74732) (3.06532,0.76392) (3.05187,0.77745) (3.03820,0.78785) (3.02438,0.79507) (3.01047,0.79909) (2.99651,0.79990) (2.98257,0.79748) (2.96869,0.79186) (2.95494,0.78304) (2.94138,0.77107) (2.92805,0.75600) (2.91500,0.73788) (2.90231,0.71680) (2.89000,0.69282) (2.87814,0.66606) (2.86677,0.63661) (2.85593,0.60460) (2.84568,0.57016) (2.83604,0.53342) (2.82707,0.49453) (2.81879,0.45365) (2.81124,0.41094) (2.80446,0.36658) (2.79846,0.32074) (2.79327,0.27362) (2.78891,0.22539) (2.78541,0.17625) (2.78276,0.12640) (2.78100,0.07604) (2.78011,0.02538) (2.78011,-0.02538) (2.78100,-0.07604) (2.78276,-0.12640) (2.78541,-0.17625) (2.78891,-0.22539) (2.79327,-0.27362) (2.79846,-0.32074) (2.80446,-0.36658) (2.81124,-0.41094) (2.81879,-0.45365) (2.82707,-0.49453) (2.83604,-0.53342) (2.84568,-0.57016) (2.85593,-0.60460) (2.86677,-0.63661) (2.87814,-0.66606) (2.89000,-0.69282) (2.90231,-0.71680) (2.91500,-0.73788) (2.92805,-0.75600) (2.94138,-0.77107) (2.95494,-0.78304) (2.96869,-0.79186) (2.98257,-0.79748) (2.99651,-0.79990) (3.01047,-0.79909) (3.02438,-0.79507) (3.03820,-0.78785) (3.05187,-0.77745) (3.06532,-0.76392) (3.07851,-0.74732) (3.09139,-0.72771) (3.10390,-0.70516) (3.11599,-0.67978) (3.12761,-0.65166) (3.13872,-0.62092) (3.14927,-0.58767) (3.15922,-0.55206) (3.16853,-0.51423) (3.17716,-0.47433) (3.18508,-0.43251) (3.19225,-0.38896) (3.19864,-0.34384) (3.20424,-0.29733) (3.20902,-0.24963) (3.21295,-0.20092) (3.21602,-0.15140) (3.21823,-0.10127) (3.21956,-0.05074) (3.22000,-0.00000)};\draw[object] (-3.00000,0.80000)--(3.00000,0.80000);\draw[object] (-3.00000,-0.80000)--(3.00000,-0.80000);\draw[guide] (-3.80000,0.00000)--(3.80000,0.00000);\draw[axis] (0.00000,-1.30000)--(0.00000,2.00000);\draw[traj] plot coordinates {(1.42000,0.00000) (1.41956,0.05074) (1.41823,0.10127) (1.41602,0.15140) (1.41295,0.20092) (1.40902,0.24963) (1.40424,0.29733) (1.39864,0.34384) (1.39225,0.38896) (1.38508,0.43251) (1.37716,0.47433) (1.36853,0.51423) (1.35922,0.55206) (1.34927,0.58767) (1.33872,0.62092) (1.32761,0.65166) (1.31599,0.67978) (1.30390,0.70516) (1.29139,0.72771) (1.27851,0.74732) (1.26532,0.76392) (1.25187,0.77745) (1.23820,0.78785) (1.22438,0.79507) (1.21047,0.79909) (1.19651,0.79990) (1.18257,0.79748) (1.16869,0.79186) (1.15494,0.78304) (1.14138,0.77107) (1.12805,0.75600) (1.11500,0.73788) (1.10231,0.71680) (1.09000,0.69282) (1.07814,0.66606) (1.06677,0.63661) (1.05593,0.60460) (1.04568,0.57016) (1.03604,0.53342) (1.02707,0.49453) (1.01879,0.45365) (1.01124,0.41094) (1.00446,0.36658) (0.99846,0.32074) (0.99327,0.27362) (0.98891,0.22539) (0.98541,0.17625) (0.98276,0.12640) (0.98100,0.07604) (0.98011,0.02538) (0.98011,-0.02538) (0.98100,-0.07604) (0.98276,-0.12640) (0.98541,-0.17625) (0.98891,-0.22539) (0.99327,-0.27362) (0.99846,-0.32074) (1.00446,-0.36658) (1.01124,-0.41094) (1.01879,-0.45365) (1.02707,-0.49453) (1.03604,-0.53342) (1.04568,-0.57016) (1.05593,-0.60460) (1.06677,-0.63661) (1.07814,-0.66606) (1.09000,-0.69282) (1.10231,-0.71680) (1.11500,-0.73788) (1.12805,-0.75600) (1.14138,-0.77107) (1.15494,-0.78304) (1.16869,-0.79186) (1.18257,-0.79748) (1.19651,-0.79990) (1.21047,-0.79909) (1.22438,-0.79507) (1.23820,-0.78785) (1.25187,-0.77745) (1.26532,-0.76392) (1.27851,-0.74732) (1.29139,-0.72771) (1.30390,-0.70516) (1.31599,-0.67978) (1.32761,-0.65166) (1.33872,-0.62092) (1.34927,-0.58767) (1.35922,-0.55206) (1.36853,-0.51423) (1.37716,-0.47433) (1.38508,-0.43251) (1.39225,-0.38896) (1.39864,-0.34384) (1.40424,-0.29733) (1.40902,-0.24963) (1.41295,-0.20092) (1.41602,-0.15140) (1.41823,-0.10127) (1.41956,-0.05074) (1.42000,-0.00000)};\draw[traj] plot coordinates {(1.62000,0.00000) (1.61956,0.05074) (1.61823,0.10127) (1.61602,0.15140) (1.61295,0.20092) (1.60902,0.24963) (1.60424,0.29733) (1.59864,0.34384) (1.59225,0.38896) (1.58508,0.43251) (1.57716,0.47433) (1.56853,0.51423) (1.55922,0.55206) (1.54927,0.58767) (1.53872,0.62092) (1.52761,0.65166) (1.51599,0.67978) (1.50390,0.70516) (1.49139,0.72771) (1.47851,0.74732) (1.46532,0.76392) (1.45187,0.77745) (1.43820,0.78785) (1.42438,0.79507) (1.41047,0.79909) (1.39651,0.79990) (1.38257,0.79748) (1.36869,0.79186) (1.35494,0.78304) (1.34138,0.77107) (1.32805,0.75600) (1.31500,0.73788) (1.30231,0.71680) (1.29000,0.69282) (1.27814,0.66606) (1.26677,0.63661) (1.25593,0.60460) (1.24568,0.57016) (1.23604,0.53342) (1.22707,0.49453) (1.21879,0.45365) (1.21124,0.41094) (1.20446,0.36658) (1.19846,0.32074) (1.19327,0.27362) (1.18891,0.22539) (1.18541,0.17625) (1.18276,0.12640) (1.18100,0.07604) (1.18011,0.02538) (1.18011,-0.02538) (1.18100,-0.07604) (1.18276,-0.12640) (1.18541,-0.17625) (1.18891,-0.22539) (1.19327,-0.27362) (1.19846,-0.32074) (1.20446,-0.36658) (1.21124,-0.41094) (1.21879,-0.45365) (1.22707,-0.49453) (1.23604,-0.53342) (1.24568,-0.57016) (1.25593,-0.60460) (1.26677,-0.63661) (1.27814,-0.66606) (1.29000,-0.69282) (1.30231,-0.71680) (1.31500,-0.73788) (1.32805,-0.75600) (1.34138,-0.77107) (1.35494,-0.78304) (1.36869,-0.79186) (1.38257,-0.79748) (1.39651,-0.79990) (1.41047,-0.79909) (1.42438,-0.79507) (1.43820,-0.78785) (1.45187,-0.77745) (1.46532,-0.76392) (1.47851,-0.74732) (1.49139,-0.72771) (1.50390,-0.70516) (1.51599,-0.67978) (1.52761,-0.65166) (1.53872,-0.62092) (1.54927,-0.58767) (1.55922,-0.55206) (1.56853,-0.51423) (1.57716,-0.47433) (1.58508,-0.43251) (1.59225,-0.38896) (1.59864,-0.34384) (1.60424,-0.29733) (1.60902,-0.24963) (1.61295,-0.20092) (1.61602,-0.15140) (1.61823,-0.10127) (1.61956,-0.05074) (1.62000,-0.00000)};\fill (0.00000,0.00000) circle(.035);\node[below left] at (0.00000,0.00000) {$C$};\node[above right] at (0.00000,2.00000) {$z$};\draw[<->,thin,color=muted] (0.00000,-1.60000)--(1.30000,-1.60000);\node[below] at (0.65000,-1.60000) {$s$};\node[right] at (1.50000,1.40000) {$\dd m=(m/L)\dd s$};\begin{scope}[shift={(8.00000,0.00000)}]
\draw[axis] (0.00000,0.00000)--(2.25000,0.00000);\draw[axis] (0.00000,-1.40000)--(0.00000,2.25000);\node[below right] at (2.25000,0.00000) {$y$};\node[above right] at (0.00000,2.25000) {$z$};\draw[object] (0.00000,0.00000) circle(1.30000);\draw[thin,color=muted] (-1.50000,0.00000)--(0.00000,0.00000);\node[below right] at (0.00000,0.00000) {$\odot x$};\node[above] at (0.00000,3.00000) {圆片直径轴惯量相同};\node[below] at (0.50000,-2.20000) {$\dd I_y=\dd I_z=R^2\dd m/4$};
\end{scope}\node[below] at (4.00000,-3.70000) {先求圆片直径惯量，再加轴向距离平方};
\end{tikzpicture}
\figcaption{解图力试2.6E\quad 圆片直径惯量加平行轴项，积分得到竖直轴惯量。}
\end{minipage}\end{center}
解图力试2.6E把圆柱切成薄圆片。圆柱绕质心铅垂轴的惯量$I_z$不是柱轴惯量。取沿柱轴坐标$s$、厚度$\dd s$的薄圆片，$\dd m=(m/L)\dd s$。圆片绕中心法线的惯量为$R^2\dd m/2$，由面内两个直径方向对称及垂直轴定理，其直径惯量为$R^2\dd m/4$。再将直径轴平移到过$C$的铅垂轴，
\[
\dd I_z=\left(\frac{R^2}{4}+s^2\right)\frac mL\dd s,
\]
\[
I_z=\int_{-L/2}^{L/2}
\left(\frac{R^2}{4}+s^2\right)\frac mL\dd s
=\frac m{12}(L^2+3R^2).
\]
积分中$R^2/4$项来自每片内部的直径惯量，$s^2$项来自片心相对质心轴的距离；两者分别积分为$mR^2/4$和$mL^2/12$。线性方程$I_z\ddot\theta+k_\theta\theta=0$给
\[
\boxed{T_3=2\pi\sqrt{\frac{4lI_z}{mgD^2}}
=2\pi\sqrt{\frac{(L^2+3R^2)l}{3gD^2}}}.
\]
质心升高按转角的平方变化，重力势能对角度求导后得到线性的恢复力矩。再沿解图力试2.6D的俯视图读张力方向。下吊点朝偏航方向横移 $D\theta/2$，每条线向回的切向分力约为 $(mg/2)(D\theta/2l)$。两支分力各以 $D/2$ 作力臂，方向都使 $\theta$ 减小，所以
\[
\tau_z \simeq -2\left(\frac{mg}{2}\frac{D\theta}{2l}\right)\frac D2
=-\frac{mgD^2}{4l}\theta.
\]
偏航的力矩也可逐项写为$M_z=-2(F\phi)(D/2)=-FD\phi$。平衡附近$2F\simeq mg$，代入$\phi=D\theta/(2l)$得到$M_z=-mgD^2\theta/(4l)$。恢复力矩的负号指向减小偏航角的方向；由能量时间导数$I_z\dot\theta\ddot\theta+mgD^2\theta\dot\theta/(4l)=0$也得到同一线性运动方程。圆柱的竖直速度为$\Delta\dot h\propto\theta\dot\theta$，其动能属振幅四阶，故线性偏航的惯性项只需保留$I_z\dot\theta^2/2$。



\end{document}
